\documentclass[11pt]{article}
\usepackage[margin=1in]{geometry}
\usepackage{amsmath}
\usepackage{graphicx}
\usepackage{url}
\usepackage{tikz-cd}
\usepackage{placeins}
\usepackage{float}
\usepackage{pdflscape}
\usepackage{rotating}

\setlength{\parindent}{0pt}
\setlength{\parskip}{1\baselineskip}

\begin{document}

\title{Modelo de equilibrio general computable de doce sectores para la evaluación de políticas tributarias}
\author{Departamento de Estudios Económicos y Tributarios, SII}
\date{Septiembre 2026}
\maketitle

\section*{Resumen ejecutivo}

Este documento desarrolla un modelo de equilibrio general computable estático y determinístico de doce sectores para Chile, calibrado con una matriz de contabilidad social de 2023. Su objetivo es evaluar los efectos macroeconómicos y recaudatorios de cambios tributarios específicos. El modelo incorpora un hogar representativo, doce sectores productivos, trece bienes, un gobierno, firmas, un sector externo y un bloque de ahorro--inversión. Asimismo, utiliza mecanismos CET y Armington para representar la asignación entre ventas domésticas, exportaciones e importaciones. En los mercados de factores, el modelo admite capacidad ociosa del capital ---con una utilización base anclada en 82\% de la capacidad instalada--- y una fuerza de trabajo que responde al empleo. Ambos márgenes se calibran para que la respuesta del desempleo al producto sea coherente con la evidencia disponible para Chile: el coeficiente de Okun implícito, medido respecto del valor agregado real, se sitúa entre 0.276 y 0.327 según el choque. Los impuestos indirectos se representan como cuñas de precio entre el precio básico y el precio de comprador, y no como usos de bienes: el IVA grava la demanda final y otorga crédito pleno al uso intermedio, mientras que los impuestos específicos sobre productos se modelan como montos fijos por unidad física que paga únicamente la demanda final.

Los resultados muestran que las cuatro políticas tributarias evaluadas generan efectos agregados diferenciados y operan mediante distintos mecanismos de transmisión. Un aumento de un punto porcentual en la tasa del IVA ---equivalente a pasar de 19\% a 20\%--- reduce el valor agregado real a costo de factores en 0.9105\% (el PIB real medido por el gasto cae 0.9455\%), el consumo real en 1.2603\% y la inversión real en 0.6549\%. La tasa de desempleo aumenta entre 0.2879 y 0.3313 puntos porcentuales, según el tipo de trabajo, y la utilización del capital cae de 82.00\% a 81.01\% de la capacidad instalada. El efecto conjunto de tasa y base asciende a 0.2371 puntos del PIB inicial: la contracción de las bases absorbe 47\% del rendimiento mecánico de 0.4503 puntos. La variación nominal observada de la recaudación total es de 0.1755 puntos del PIB, debido a un efecto precio de -0.0616 puntos.

Un aumento de 10\% en el impuesto a la renta de personas reduce el PIB real en 0.4644\%, el valor agregado real en 0.4559\%, el consumo real en 0.6335\% y la inversión real en 0.2541\%. La tasa de desempleo se incrementa entre 0.1479 y 0.1577 puntos porcentuales, según el tipo de trabajo, y la utilización del capital cae de 82.00\% a 81.50\% de la capacidad instalada. El efecto conjunto de tasa y base eleva la recaudación en 0.0837 puntos del PIB inicial, mientras que la variación nominal observada es de 0.0527 puntos, debido a un efecto precio de -0.0310 puntos.

Un alza de 10\% en los impuestos específicos sobre productos reduce el valor agregado real en 0.2166\% (el PIB real por el gasto cae 0.2289\%), el consumo real en 0.3037\% y la inversión real en 0.1615\%. Su rendimiento mecánico es de 0.0970 puntos del PIB y la erosión de la propia base absorbe apenas 5\% de ese monto; la contracción de las demás bases, en cambio, resta 0.0503 puntos, de modo que el efecto conjunto de tasa y base sobre la recaudación total es de 0.0419 puntos y la variación nominal observada, de 0.0441 puntos.

En contraste, una disminución de un punto porcentual en la tasa del impuesto de primera categoría (IDPC) eleva el PIB real en 0.1733\%, el consumo real en 0.1369\% y la inversión real en 0.4116\%. También reduce la tasa de desempleo entre 0.0533 y 0.0696 puntos porcentuales, según el tipo de trabajo, aunque deteriora el saldo comercial. El efecto conjunto de tasa y base reduce la recaudación en 0.1255 puntos del PIB inicial: la expansión de las bases compensa cerca de 22\% del efecto mecánico. La variación nominal observada es de -0.1140 puntos, debido a un efecto precio positivo de 0.0115 puntos.


\vspace{0.5\baselineskip}
\noindent\textbf{Nota de versión.} Esta versión reespecifica el IVA y los impuestos específicos como cuñas de precio (Sección~\ref{sec:productos}) y ancla la utilización base del capital en 82\% de la capacidad instalada (Sección~\ref{sec:equilibrio}). Las cuatro simulaciones reportadas corresponden a la especificación vigente. Los resultados del incremento del IVA difieren de manera sustantiva de los publicados en versiones anteriores del documento, por las razones que se detallan en la Sección~\ref{sec:sim-iva}.

\section{Introducción}

Este documento presenta un modelo de equilibrio general computable (MEGC) estático y determinístico, orientado a evaluar el efecto de distintas medidas tributarias.

El modelo considera:

\begin{itemize}
  \item Un hogar representativo.
  \item Doce sectores productivos: agropecuario-silvícola y pesca; minería; industria manufacturera; electricidad, gas, agua y gestión de desechos; construcción; comercio, hoteles y restaurantes; transporte, comunicaciones y servicios de información; intermediación financiera; servicios inmobiliarios y de vivienda; servicios empresariales; servicios personales; y administración pública.
  \item Trece tipos de bienes: productos agropecuarios, silvícolas y pesqueros; minerales; productos manufacturados; electricidad, gas, agua y gestión de desechos; construcción; servicios de comercio, hoteles y restaurantes; servicios de transporte, comunicaciones e información; servicios financieros; servicios inmobiliarios y de vivienda; servicios empresariales; servicios personales; servicios de la administración pública; y otros bienes y servicios.
  \item Un gobierno.
  \item Firmas.
  \item Un sector externo.
  \item Un bloque de ahorro--inversión.
  \item Un tipo de capital, con grado de utilización endógeno.
  \item Tres tipos de mano de obra (calificación baja, media y alta), con desempleo y participación laboral endógenos.
  \item Cinco instrumentos tributarios: IVA, impuesto a la renta de personas, impuesto de primera categoría, impuestos a la producción e impuestos específicos sobre productos, además de los derechos de importación. El IVA y los impuestos específicos operan como cuñas entre el precio básico y el precio de comprador.
\end{itemize}

La calibración se realiza mediante el paquete \texttt{gEcon} para R\footnote{\url{http://gecon.r-forge.r-project.org/index.html}} y una matriz de contabilidad social (SAM) construida para Chile en 2023 a partir de la matriz insumo-producto más reciente disponible. Para su elaboración se adapta la metodología de compilación matricial propuesta por \cite{venegas_2013}.

La literatura que enmarca este estudio comprende tres dimensiones. La primera corresponde a los fundamentos teóricos del equilibrio general y su aplicación empírica al análisis de políticas públicas \cite{arrow_debreu_1954,shoven_whalley_1984,shoven_whalley_1992,dixon_jorgenson_2013}. La segunda recoge la tradición de modelos aplicados calibrados con matrices de contabilidad social \cite{pyatt_round_1985,dervis_demelo_robinson_1982}, que permite representar de manera coherente los flujos entre hogares, firmas y gobierno. La tercera comprende la evidencia metodológica sobre evaluación tributaria mediante MEGC \cite{ballard_fullerton_shoven_whalley_1985}, que orienta la interpretación de los efectos sobre la recaudación, el bienestar y la composición de la demanda agregada.

El resto del documento se organiza de la siguiente manera. Primero, se presenta la matriz de contabilidad social utilizada para calibrar el modelo y se describen las principales cuentas de la economía. Luego, se desarrollan los distintos bloques del MEGC ---hogares, producción, firmas, gobierno, sector externo y ahorro--inversión---, junto con el cierre macroeconómico adoptado. A continuación, se explica la metodología empleada para evaluar el impacto de las medidas tributarias simuladas. Posteriormente, se presentan los resultados de las simulaciones. Por último, se sintetizan los principales hallazgos y se plantean posibles extensiones del ejercicio.

\section{Matriz de Contabilidad Social (SAM) de Chile}

Esta sección presenta los datos utilizados para calibrar el modelo, que representan la economía chilena en 2023. En la matriz, cada columna registra los gastos de una cuenta y cada fila, sus ingresos. En equilibrio, la suma de cada columna debe coincidir con la de la fila correspondiente.

De acuerdo con la clasificación del Banco Central, la economía comprende trece bienes: productos agropecuarios, silvícolas y pesqueros; minerales; productos manufacturados; electricidad, gas, agua y gestión de desechos; construcción; servicios de comercio, hoteles y restaurantes; servicios de transporte, comunicaciones e información; servicios financieros; servicios inmobiliarios y de vivienda; servicios empresariales; servicios personales; servicios de la administración pública; y otros bienes y servicios. Estos bienes son producidos por doce sectores: agropecuario-silvícola y pesca; minería; industria manufacturera; electricidad, gas, agua y gestión de desechos; construcción; comercio, hoteles y restaurantes; transporte, comunicaciones y servicios de información; intermediación financiera; servicios inmobiliarios y de vivienda; servicios empresariales; servicios personales; y administración pública. Cada sector utiliza trabajo, capital e insumos intermedios, además de pagar impuestos a la producción.

La producción sectorial se obtuvo de la matriz de producción a precios básicos, incorporando el margen de comercio en la actividad de servicios. Los insumos se construyeron a partir del cuadro de utilización intermedia total a precios de usuario, mientras que el trabajo, el capital y los impuestos a la producción se obtuvieron del cuadrante de valor agregado.

El trabajo se clasificó en tres niveles de calificación ---baja, media y alta--- a partir de la distribución de remuneraciones por sector económico de la encuesta CASEN 2022. La categoría de baja calificación comprende a las personas con educación básica completa o inferior; la de calificación media, a quienes alcanzaron como máximo la educación media completa; y la de alta calificación, a quienes poseen educación superior completa o incompleta.

La oferta disponible de cada bien se compone de la producción doméstica, las importaciones provenientes del resto del mundo (ROW), el IVA y los aranceles. Esta oferta se destina a demanda intermedia, consumo de los hogares y del gobierno, exportaciones e inversión. Esta última se desagrega entre inversión propiamente tal (S-I) y variación de existencias.

Los impuestos específicos, el IVA, los aranceles y las importaciones se obtuvieron del cuadro de oferta total. Por su parte, los insumos intermedios se tomaron del cuadro de utilización intermedia total a precios de usuario, mientras que el consumo de los hogares y del gobierno, las exportaciones, la inversión y la variación de existencias se obtuvieron del cuadro de utilización final total a precios de usuario.

Los hogares obtienen su ingreso disponible a partir de las remuneraciones de los tres tipos de trabajo, las rentas del capital, las transferencias provenientes de las empresas y del gobierno, y las prestaciones sociales. Este ingreso se destina al consumo, el ahorro, el pago de contribuciones sociales y el impuesto a la renta.

La totalidad de los ingresos laborales se asigna a los hogares. Los ingresos de capital corresponden a las rentas de la producción de los hogares, netas de las remuneraciones al trabajo. Las transferencias gubernamentales provienen de las transferencias sociales en especie; las prestaciones sociales corresponden a las efectuadas por empresas y gobierno; y las transferencias empresariales se determinan residualmente para cuadrar la cuenta. El gasto en bienes se obtiene del cuadro de utilización final total a precios de usuario. Las contribuciones sociales se aproximan mediante el ahorro previsional; el impuesto a la renta, mediante los impuestos netos de subvenciones; y el ahorro, mediante la formación bruta de capital.

Las empresas perciben rentas del capital y contribuciones sociales, que destinan a transferencias a los hogares, prestaciones sociales, ahorro y pago de impuestos a la renta. Cabe destacar que el impuesto corporativo se modela como un gravamen sobre la renta del capital a nivel institucional (firmas), y no como una cuña directa sobre el costo de uso del capital en las decisiones productivas sectoriales. Por tanto, sus efectos se transmiten principalmente a través del ahorro empresarial y la inversión, más que mediante cambios directos en la demanda de factores.

Los retornos al capital de las empresas se obtienen de las rentas de la producción y de la propiedad neta registradas en las tablas de operaciones no financieras de las sociedades financieras y no financieras, más un ajuste destinado a cuadrar las cuentas. Los dividendos distribuidos a los hogares y al gobierno también se determinan mediante un ajuste. Las prestaciones sociales se aproximan mediante la variación neta de los fondos de pensiones registrada en las operaciones no financieras, neta de impuestos sobre subvenciones. Por su parte, el ahorro empresarial se obtiene de la formación bruta de capital de las empresas financieras y no financieras.

El gobierno gasta en bienes y servicios, así como en ahorro, transferencias a los hogares y prestaciones sociales. Sus ingresos provienen de las rentas del capital, las contribuciones sociales, los impuestos a la renta, el IVA, los impuestos a la producción, los impuestos específicos sobre productos y los aranceles. En la clasificación utilizada para Chile, los impuestos a la producción agrupan tributos asociados al proceso productivo y al uso de activos o permisos ---como timbres y estampillas, contribuciones de bienes raíces, permisos de circulación, patentes y derechos municipales y no municipales, e impuestos a los juegos de azar, casinos e hipódromos, entre otros---. En cambio, los impuestos específicos sobre productos corresponden a gravámenes aplicados por unidad de bien o como porcentaje de su valor de venta. Dentro de estos últimos destacan los impuestos de la Ley de Alcoholes y los impuestos específicos a combustibles, tabaco y fuentes móviles.

El resto del mundo percibe ingresos por las importaciones, las remuneraciones del trabajo extranjero y las rentas del capital que posee. Estos recursos se destinan al pago de las exportaciones de la economía, al ahorro externo y a las remuneraciones del capital doméstico situado en el extranjero. En esta representación, el ahorro externo corresponde al financiamiento neto proveniente del resto del mundo y es consistente con la diferencia entre la inversión y el ahorro doméstico. Su registro en la matriz sigue la convención contable de la SAM y puede reflejarse indirectamente en las cuentas de capital y propiedad.



La Tabla~\ref{tab:sam-resumida-2023} presenta una versión resumida de la SAM utilizada en la calibración, agregada a un sector productivo, un producto y un tipo de trabajo. Las celdas se reportan con un decimal y las posiciones sin flujo observado se muestran como 0.0.

\begin{landscape}
\begin{table}[H]
\centering
\caption{Matriz de Contabilidad Social resumida de Chile, 2023}
\label{tab:sam-resumida-2023}
\scriptsize
\setlength{\tabcolsep}{2pt}
\renewcommand{\arraystretch}{1.05}
\resizebox{\linewidth}{!}{%
\begin{tabular}{|l|r|r|r|r|r|r|r|r|r|r|r|r|r|r|r|r|r|r|}
  \hline
  Cuenta & Act & prod & Trabajo & Capital & Hogares & Empresas & Gobierno & Cont\_soc & Pres\_soc & Imp\_dir & Imp\_prod & Imptos\_Espec & IVA & Arancel & row & S-I & Var\_Exis & Total \\
  \hline
  Act & 0.0 & 498537.6 & 0.0 & 0.0 & 0.0 & 0.0 & 0.0 & 0.0 & 0.0 & 0.0 & 0.0 & 0.0 & 0.0 & 0.0 & 0.0 & 0.0 & 0.0 & 498537.6 \\
  \hline
  prod & 245058.5 & 0.0 & 0.0 & 0.0 & 170627.0 & 0.0 & 43183.6 & 0.0 & 0.0 & 0.0 & 0.0 & 0.0 & 0.0 & 0.0 & 86815.4 & 68824.6 & -3002.6 & 611506.5 \\
  \hline
  Trabajo & 114601.6 & 0.0 & 0.0 & 0.0 & 0.0 & 0.0 & 0.0 & 0.0 & 0.0 & 0.0 & 0.0 & 0.0 & 0.0 & 0.0 & 0.0 & 0.0 & 0.0 & 114601.6 \\
  \hline
  Capital & 134717.9 & 0.0 & 0.0 & 0.0 & 0.0 & 0.0 & 0.0 & 0.0 & 0.0 & 0.0 & 0.0 & 0.0 & 0.0 & 0.0 & 0.0 & 0.0 & 0.0 & 134717.9 \\
  \hline
  Hogares & 0.0 & 0.0 & 114601.6 & 29078.5 & 0.0 & 34181.8 & 2292.6 & 0.0 & 20407.9 & 0.0 & 0.0 & 0.0 & 0.0 & 0.0 & 0.0 & 0.0 & 0.0 & 200562.4 \\
  \hline
  Empresas & 0.0 & 0.0 & 0.0 & 102970.9 & 0.0 & 0.0 & 0.0 & 15585.8 & 0.0 & 0.0 & 0.0 & 0.0 & 0.0 & 0.0 & 0.0 & 0.0 & 0.0 & 118556.7 \\
  \hline
  Gobierno & 0.0 & 0.0 & 0.0 & 2668.6 & 0.0 & 8772.7 & 0.0 & 0.0 & 0.0 & 18540.9 & 4159.7 & 2743.9 & 24332.0 & 1541.6 & 0.0 & 0.0 & 0.0 & 62759.4 \\
  \hline
  Cont\_soc & 0.0 & 0.0 & 0.0 & 0.0 & 15585.8 & 0.0 & 0.0 & 0.0 & 0.0 & 0.0 & 0.0 & 0.0 & 0.0 & 0.0 & 0.0 & 0.0 & 0.0 & 15585.8 \\
  \hline
  Pres\_soc & 0.0 & 0.0 & 0.0 & 0.0 & 0.0 & 9387.2 & 11020.7 & 0.0 & 0.0 & 0.0 & 0.0 & 0.0 & 0.0 & 0.0 & 0.0 & 0.0 & 0.0 & 20407.9 \\
  \hline
  Imp\_dir & 0.0 & 0.0 & 0.0 & 0.0 & 5106.7 & 13434.2 & 0.0 & 0.0 & 0.0 & 0.0 & 0.0 & 0.0 & 0.0 & 0.0 & 0.0 & 0.0 & 0.0 & 18540.9 \\
  \hline
  Imp\_prod & 4159.7 & 0.0 & 0.0 & 0.0 & 0.0 & 0.0 & 0.0 & 0.0 & 0.0 & 0.0 & 0.0 & 0.0 & 0.0 & 0.0 & 0.0 & 0.0 & 0.0 & 4159.7 \\
  \hline
  Imptos\_Espec & 0.0 & 2743.9 & 0.0 & 0.0 & 0.0 & 0.0 & 0.0 & 0.0 & 0.0 & 0.0 & 0.0 & 0.0 & 0.0 & 0.0 & 0.0 & 0.0 & 0.0 & 2743.9 \\
  \hline
  IVA & 0.0 & 24332.0 & 0.0 & 0.0 & 0.0 & 0.0 & 0.0 & 0.0 & 0.0 & 0.0 & 0.0 & 0.0 & 0.0 & 0.0 & 0.0 & 0.0 & 0.0 & 24332.0 \\
  \hline
  Arancel & 0.0 & 1541.6 & 0.0 & 0.0 & 0.0 & 0.0 & 0.0 & 0.0 & 0.0 & 0.0 & 0.0 & 0.0 & 0.0 & 0.0 & 0.0 & 0.0 & 0.0 & 1541.6 \\
  \hline
  row & 0.0 & 84351.4 & 0.0 & 0.0 & 0.0 & 2464.1 & 0.0 & 0.0 & 0.0 & 0.0 & 0.0 & 0.0 & 0.0 & 0.0 & 0.0 & 0.0 & 0.0 & 86815.4 \\
  \hline
  S-I & 0.0 & 0.0 & 0.0 & 0.0 & 9242.8 & 50316.8 & 6262.4 & 0.0 & 0.0 & 0.0 & 0.0 & 0.0 & 0.0 & 0.0 & 0.0 & 0.0 & 0.0 & 65822.0 \\
  \hline
  Var\_Exis & 0.0 & 0.0 & 0.0 & 0.0 & 0.0 & 0.0 & 0.0 & 0.0 & 0.0 & 0.0 & 0.0 & 0.0 & 0.0 & 0.0 & 0.0 & -3002.6 & 0.0 & -3002.6 \\
  \hline
  Total & 498537.6 & 611506.5 & 114601.6 & 134717.9 & 200562.4 & 118556.7 & 62759.4 & 15585.8 & 20407.9 & 18540.9 & 4159.7 & 2743.9 & 24332.0 & 1541.6 & 86815.4 & 65822.0 & -3002.6 & 0.0 \\
  \hline
\end{tabular}%
}
\end{table}
\end{landscape}
\FloatBarrier

\section{Componentes del modelo}

Esta sección describe los bloques que conforman la economía del modelo. La producción de los doce sectores se distribuye entre ventas domésticas y exportaciones mediante una función CET, mientras que la absorción interna se construye mediante un agregado Armington de bienes domésticos e importados. Los conjuntos de sectores, productos y bienes consumidos por el hogar se definen a continuación:
\[
\begin{aligned}
\mathrm{SEC}=\{&\mathrm{ActAgri},\mathrm{ActMin},\mathrm{ActManu},\mathrm{ActElectAguaGas}, \\
&\mathrm{ActConst},\mathrm{ActComHotRest},\mathrm{ActTransComuInf},\mathrm{ActIntFin}, \\
&\mathrm{ActServEmp},\mathrm{ActServPers},\mathrm{ActAdmPub}\}.
\end{aligned}
\]

\[
\begin{aligned}
\mathrm{PRODS}=\{&\mathrm{ProdAgri},\mathrm{ProdMin},\mathrm{ProdManu},\mathrm{ProdElectAguaGas}, \\
&\mathrm{ProdConst},\mathrm{ProdComHotRest},\mathrm{ProdTransComuInf},\mathrm{ProdIntFin}, \\
&\mathrm{ProdInmoViv},\mathrm{ProdServEmp},\mathrm{ProdServPers},\mathrm{ProdAdmPub}, \\
&\mathrm{ProdOtros}\}.
\end{aligned}
\]

\[
\begin{aligned}
\mathrm{PCON}=\{&\mathrm{ProdAgri},\mathrm{ProdManu},\mathrm{ProdElectAguaGas},\mathrm{ProdComHotRest}, \\
&\mathrm{ProdTransComuInf},\mathrm{ProdIntFin},\mathrm{ProdInmoViv},\mathrm{ProdServEmp}, \\
&\mathrm{ProdServPers},\mathrm{ProdAdmPub},\mathrm{ProdOtros}\}.
\end{aligned}
\]

La economía también incluye un hogar representativo, un gobierno, un sector externo y un bloque de ahorro--inversión.

\subsection{Bloque del consumidor}

El consumidor corresponde a un hogar representativo que maximiza su utilidad mediante el consumo de los bienes compuestos incluidos en \(\mathrm{PCON}\). En la implementación, las demandas de productos mineros y de construcción se fijan en niveles residuales exógenos, por lo que estos bienes no ingresan directamente en la función de utilidad:

\[
\max_{\left(D^{p}\right)_{p \in \mathrm{PCON}}} \;
U = \left(\sum_{p \in \mathrm{PCON}} \alpha^{p} \left(D^{p}\right)^{\frac{\omega - 1}{\omega}}\right)^{\frac{\omega}{\omega - 1}}
\]

Esta optimización está sujeta a las siguientes restricciones:

\[
INC = \sum_{p \in \mathrm{PCON}} pq^{p} D^{p} + SAV + \mathrm{CONT}^{\mathrm{SOC}}
\]

\[
INC = BTINC - PITAX
\]

\[
BTINC = L^{h} w + K^{h} p^{k} + \mathrm{TR}^{\mathrm{EMP}} + \mathrm{TR}^{\mathrm{GOV}} + PS^{h} + FACTROW^{K} + FACTROW^{L}
\]

\[
PITAX = pit \cdot \left(K^{h} p^{k} + L^{h} w + \mathrm{TR}^{\mathrm{GOV}} + \mathrm{TR}^{\mathrm{EMP}} + PS^{h} + FACTROW^{K} + FACTROW^{L}\right)
\]

En estas expresiones, $D^{p}$ es la demanda del hogar por el bien compuesto $p$, $\omega$ se fija en 1.5 y $\alpha^{p}$ representa el parámetro de participación del bien en el consumo, calibrado para reproducir la SAM. El precio relevante para los hogares es $pq^{p}$, el precio de comprador del bien compuesto, que incorpora el impuesto específico por unidad física y el IVA sobre la demanda final (Sección~\ref{sec:productos}). En consecuencia, la sustitución entre bienes en la función de utilidad responde a los precios efectivamente pagados por el hogar y no al precio básico $p^{p}$. Asimismo, $SAV$ representa el ahorro de los hogares y $CONT^{SOC}$, las contribuciones sociales que estos pagan.

Por su parte, $INC$ es el ingreso disponible de los hogares, $BTINC$ es el ingreso antes de impuestos y $pit$ es la tasa efectiva del impuesto a la renta de personas. En la descomposición de $BTINC$, $L^{h}$ representa la dotación de trabajo; $w$, el salario; $K^{h}$, la dotación de capital de los hogares; $p^{k}$, el retorno del capital; $TR^{EMP}$ y $TR^{GOV}$, las transferencias provenientes de las empresas y del gobierno; y $PS^{h}$, las prestaciones sociales recibidas por los hogares. Los términos $FACTROW^{K}$ y $FACTROW^{L}$ representan, respectivamente, los ingresos de capital y de trabajo provenientes del resto del mundo. De acuerdo con la SAM utilizada, ambos se fijan en cero: $FACTROW^{K}=0$ y $FACTROW^{L}=0$.

En esta formulación, el vínculo con el sector externo no ingresa en el presupuesto del hogar mediante rentas residuales del comercio, sino que se internaliza en los bloques CET y Armington a través de los precios $pe^{p}$, $pd^{p}$, $pm^{p}$ y $p^{p}$.

\subsection{Bloque de sectores}

El bloque de sectores constituye el núcleo productivo del MEGC: en cada actividad $s \in \mathrm{SEC}$, una firma representativa elige producción, factores e insumos para maximizar beneficios.

Debido a los ceros estructurales de la SAM, la especificación sectorial se divide según los insumos intermedios utilizados. Esta partición evita problemas numéricos al calibrar demandas de insumos cuyos flujos son exactamente nulos en la matriz.

\textbf{1) Variables de decisión.} Para cada sector, la firma determina \(V^{s}\) (nivel de actividad), \(K^{s}\) (capital) y la demanda de cada tipo de trabajo \(L^{l,s}\), con \(l \in \mathcal{L}=\{1,2,3\}\). La demanda de insumos intermedios \(X^{p,s}\) no se elige de manera independiente, sino que queda determinada por coeficientes fijos una vez definido \(V^{s}\).

\textbf{2) Función objetivo.} El beneficio sectorial se representa como:

\[
\pi^{s} = V^{s}\left(\sum_{p \in \mathrm{PRODS}} pz^{p}\, make\_share^{p,s}\right)\left(1-tps^{s}-tes^{s}\right)
- K^{s} p^{k} - \sum_{l \in \mathcal{L}} w^{l} L^{l,s} - \sum_{p \in \mathcal{I}(s)} io^{p,s} V^{s} pint^{p}
\]

donde \(\mathcal{I}(s)\) es el conjunto de insumos efectivamente utilizados por cada grupo sectorial e \(io^{p,s}\) es el requerimiento fijo del insumo compuesto $p$ por unidad de actividad. El primer término valora la producción al precio del productor $pz^{p}$, de manera consistente con el bloque CET. A este ingreso se restan los costos del capital, el trabajo y los insumos intermedios, valorados al precio de insumo $pint^{p}$. Este precio excluye los impuestos específicos y, salvo el residuo de calibración descrito más adelante, también el IVA, de modo que el crédito fiscal impide que el impuesto se acumule en cascada a lo largo de la cadena productiva \cite{ebrill_etal_2001}.

\textbf{3) Tecnología de producción.} La actividad productiva combina capital y los tres tipos de trabajo mediante una Cobb--Douglas de valor agregado:

\[
V^{s} = \gamma^{s} \left(K^{s}\right)^{\beta_{K}^{s}} \prod_{l \in \mathcal{L}} \left(L^{l,s}\right)^{\beta_{L,l}^{s}}.
\]

Los insumos intermedios, en cambio, siguen una estructura Leontief con coeficientes fijos proporcionales al nivel de actividad:

\[
X^{p,s} = io^{p,s} V^{s}, \qquad p \in \mathcal{I}(s)
\]

Así, el margen de sustitución dentro del bloque productivo opera entre capital y los distintos tipos de trabajo, mientras que la canasta intermedia requerida por cada sector mantiene las proporciones observadas en la SAM.

\textbf{4) Identidades contables y matriz de producción.}

\[
YAP^{p,s} = V^{s} \cdot make\_share^{p,s}
\]

\[
revenue^{s} = \sum_{p \in \mathrm{PRODS}} YAP^{p,s} pz^{p}
\]

\[
net\_revenue^{s} = revenue^{s} \left(1-tps^{s}-tes^{s}\right)
\]

\[
va\_nom^{s} = net\_revenue^{s} - \sum_{p \in \mathcal{I}(s)} io^{p,s} V^{s} p^{p}
\]

La identidad \(YAP^{p,s}\) distribuye la producción sectorial entre los bienes del productor de acuerdo con la matriz \(make\_share\) y las proporciones observadas en la SAM. A partir de esta asignación se calculan el ingreso bruto, el ingreso neto de impuestos operativos y, una vez descontado el costo de los insumos intermedios, el valor agregado nominal \(va\_nom^{s}\). Posteriormente, la oferta por producto ingresa en el bloque CET, que determina su distribución entre ventas domésticas y exportaciones.

\textbf{5) Calibración.} El estado base reproduce exactamente la SAM y permite inferir los parámetros tecnológicos.\footnote{Durante la calibración se impone \(V^{s}=vst^{s}\Rightarrow \gamma^{s}\), \(K^{s}p^{k}=va\_nom^{s}\beta_{K}^{s}\Rightarrow \beta_{K}^{s}\) y \(w^{l}L^{l,s}=va\_nom^{s}\beta_{L,l}^{s}\Rightarrow \beta_{L,l}^{s}\) para cada tipo de trabajo \(l\). Las demandas intermedias se recuperan mediante los coeficientes fijos \(io^{p,s}\) observados en la SAM.}

En síntesis, cada sector resuelve un problema de optimización microeconómica. Ante choques de política o de demanda, ajusta el empleo y el capital dentro de una tecnología Cobb--Douglas de valor agregado, mientras que la demanda intermedia varía en proporciones fijas consistentes con la SAM.

\subsection{Bloque de productos y comercio exterior}
\label{sec:productos}

El bloque de productos conecta la producción con el comercio exterior y la absorción doméstica. A diferencia de una versión simplificada del modelo, la producción interna y las importaciones no se igualan directamente con la demanda final. Cada producto pasa por dos etapas: primero, una transformación CET entre ventas domésticas y exportaciones; luego, una agregación Armington entre oferta doméstica e importaciones para formar el bien compuesto utilizado internamente.

En la implementación del archivo \texttt{.gcn}, las ecuaciones se agrupan por familias de productos para tratar los ceros estructurales de la SAM, de forma análoga a la partición técnica de los sectores. Con ello se evita sumar flujos que son exactamente nulos en el escenario base.

\textbf{1) Transformación CET entre ventas domésticas y exportaciones.} Siguiendo la especificación de elasticidad de transformación constante \cite{powell_gruen_1968}, para los bienes exportables el productor resuelve:

\[
\max_{Z^{p},DS^{p},EX^{p}} \; \Pi_{CET}^{p}=pe^{p}EX^{p}+pd^{p}DS^{p}-pz^{p}Z^{p}
\]

bajo la restricción tecnológica:

\[
Z^{p}=a_{t}^{p}\left[\gamma_{e}^{p}\left(EX^{p}\right)^{\rho_{t}}+\left(1-\gamma_{e}^{p}\right)\left(DS^{p}\right)^{\rho_{t}}\right]^{1/\rho_{t}},
\qquad \rho_{t}=\frac{\sigma_{t}+1}{\sigma_{t}}
\]

En esta expresión, $Z^{p}$ es la oferta total del bien del productor; $DS^{p}$, las ventas domésticas; $EX^{p}$, las exportaciones; $pz^{p}$, el precio del productor; $pd^{p}$, el precio de venta interna; y $pe^{p}$, el precio de exportación en moneda doméstica. Para los bienes no exportables, se impone $DS^{p}=Z^{p}$, $EX^{p}=0$ y $pz^{p}=pd^{p}$.

\textbf{2) Agregado Armington para la absorción doméstica.} El bien compuesto usado por hogares, gobierno, inversión e insumos intermedios supone diferenciación imperfecta por origen \cite{armington_1969} y se obtiene resolviendo:

\[
\max_{Q^{p},DD^{p},M^{p}} \; \Pi_{ARM}^{p}=p^{p}Q^{p}-pd^{p}DD^{p}-pm^{p}M^{p}
\]

bajo la restricción:

\[
Q^{p}=a_{q}^{p}\left[\delta_{m}^{p}\left(M^{p}\right)^{\rho_{q}}+\left(1-\delta_{m}^{p}\right)\left(DD^{p}\right)^{\rho_{q}}\right]^{1/\rho_{q}},
\qquad \rho_{q}=\frac{\sigma_{q}-1}{\sigma_{q}}
\]

con $DD^{p}=DS^{p}$; es decir, el componente doméstico del agregado Armington coincide con las ventas internas determinadas por CET. Aquí, $Q^{p}$ es el bien compuesto; $M^{p}$, las importaciones; $p^{p}$, el precio compuesto; y $pm^{p}$, el precio de importación con arancel. En la calibración base, los parámetros $\delta_{m}^{p}$ y $\gamma_{e}^{p}$ se obtienen directamente de las participaciones observadas en la SAM.

\textbf{3) Absorción doméstica y vaciado del mercado compuesto.} La absorción doméstica se separa según el tipo de comprador, porque el tratamiento tributario difiere entre demanda final y uso intermedio. En unidades físicas:

\[
FD^{p}=D^{p}+GG^{p}+I^{data,p}\,inv\_adj,
\qquad
XI^{p}=\sum_{s\in\mathrm{SEC}}X^{p,s}
\]

\[
TD\_Dom^{p}=FD^{p}+VEXIST^{p}+XI^{p}
\]

donde $FD^{p}$ es la demanda final gravada ---consumo de los hogares $D^{p}$, gasto del gobierno $GG^{p}$ e inversión endógena $I^{data,p}\,inv\_adj$---; $XI^{p}$, la demanda intermedia agregada; y $VEXIST^{p}$, la variación de existencias. El vaciado del mercado del bien compuesto es entonces una igualdad puramente física:

\[
Q^{p}=TD\_Dom^{p}
\]

Esta formulación reemplaza la utilizada en versiones anteriores del modelo, en la que los impuestos indirectos se descontaban de la oferta como si absorbieran unidades físicas, mediante $Q^{p}+VAT^{p}/p^{p}+TESP^{p}/p^{p}=TD\_Dom^{p}$. En el sistema de cuentas nacionales, los impuestos sobre los productos son un componente de valoración que separa el precio básico del precio de comprador y no un uso de bienes \cite{sna_2008}. La especificación vigente respeta esa distinción: los impuestos indirectos desplazan precios y no cantidades, de manera análoga al tratamiento estándar de los impuestos sobre ventas en los MEGC aplicados \cite{lofgren_harris_robinson_2002}.

\textbf{4) Impuestos indirectos como cuñas de precio.} Cada producto tiene dos precios de usuario, según quién lo compre:

\[
pq^{p}=\left(p^{p}+\tau_{e}^{p}\right)\left(1+t_{iva}^{p}\right),
\qquad
pint^{p}=p^{p}\left(1+t_{iva,int}^{p}\right)
\]

El precio $pq^{p}$ es el precio de comprador que enfrentan los hogares, el gobierno y la inversión. Incorpora primero el impuesto específico $\tau_{e}^{p}$, expresado como un monto fijo por unidad física, y luego el IVA, que se aplica sobre el precio ya gravado con el específico. Este orden reproduce la regla del artículo 15 del DL 825, según la cual la base imponible del IVA incluye los demás tributos que afectan a la operación, con excepción del propio IVA \cite{dl825}. El precio $pint^{p}$ es el precio de los insumos intermedios: no incluye el impuesto específico ---se supone que el consumo intermedio no lo paga--- y su tasa $t_{iva,int}^{p}$ es nula salvo por el residuo de calibración que se describe en el punto 6.

La recaudación de ambos instrumentos se obtiene aplicando esas cuñas a las bases físicas correspondientes:

\[
VAT^{p}=t_{iva}^{p}\left(p^{p}+\tau_{e}^{p}\right)FD^{p}+t_{iva,int}^{p}\,p^{p}\,XI^{p},
\qquad
TESP^{p}=\tau_{e}^{p}\,FD^{p}
\]

La segunda identidad tiene una implicancia relevante para la evaluación fiscal: como el impuesto específico es un monto por unidad física, su recaudación sigue a los volúmenes y es invariante a los precios. Un impuesto unitario y uno ad valorem equivalente en la base recaudan lo mismo, pero responden de manera distinta a cualquier perturbación que altere precios relativos o el nivel general de precios, diferencia bien documentada en la literatura sobre diseño de impuestos selectivos \cite{keen_1998,delipalla_keen_1992,cnossen_2005}. Esta propiedad también implica que el valor real de la recaudación específica se erosiona con la inflación cuando el gravamen no se indexa; en Chile, los impuestos a los combustibles y al tabaco están expresados en UTM, lo que provee esa indexación \cite{ley18502}.

\textbf{5) Recaudación aduanera y vínculo con precios externos.} El bloque completa las identidades tributarias y de precios de frontera:

\[
ARANCEL^{p}=tar\_rate^{p}\,pworldm^{p}\,exr\,M^{p}
\]

\[
pm^{p}=pworldm^{p}\,exr\left(1+tar\_rate^{p}\right),
\qquad
pe^{p}=pworlde^{p}\,exr
\]

donde $tar\_rate^{p}$ es la tasa arancelaria efectiva por producto. En conjunto, el bloque determina la relación entre el precio del productor $pz^{p}$, el precio doméstico $pd^{p}$, el precio importado $pm^{p}$, el precio básico del compuesto $p^{p}$ y los dos precios de usuario $pq^{p}$ y $pint^{p}$. En la versión actual del cierre externo, el tipo de cambio $exr$ es endógeno, de modo que estos precios transmiten el ajuste necesario para compatibilizar el equilibrio comercial con un nivel fijo de ahorro externo.

\textbf{6) Calibración de las cuñas tributarias.} La SAM registra la recaudación de IVA y de impuestos específicos por producto, pero no distingue qué parte recae sobre la demanda final. La calibración asigna toda la fila de IVA a la demanda final gravada, con un tope igual a la tasa legal $\bar{t}=0.19$, y envía el excedente al uso intermedio como residuo transitorio:

\[
IVA^{p}_{FD}=\min\left\{IVA^{p},\;\frac{\bar{t}}{1+\bar{t}}\,FD^{p}_{c}\right\},
\qquad
t_{iva}^{p}=\frac{IVA^{p}_{FD}}{FD^{p}_{c}-IVA^{p}_{FD}},
\qquad
t_{iva,int}^{p}=\frac{IVA^{p}-IVA^{p}_{FD}}{XI^{p}_{c}-\left(IVA^{p}-IVA^{p}_{FD}\right)}
\]

donde el subíndice $c$ indica valoración a precios de comprador de la SAM. El residuo $t_{iva,int}^{p}$ no es un artificio contable sin contenido económico: recoge el IVA que no se recupera como crédito porque la cadena de débitos y créditos se interrumpe en actividades exentas, situación que genera efectos de cascada sobre los precios de los insumos \cite{ebrill_etal_2001,keen_2013}. En una etapa posterior este residuo se reemplazará por estimaciones de IVA no recuperable por sector construidas a partir del formulario 29.

El impuesto específico unitario se recupera de la condición de que la recaudación observada coincida con $\tau_{e}^{p}$ veces la demanda final física, valorada al precio básico del año base:

\[
\tau_{e}^{p}=\frac{p^{p,b}\;TESP^{p}}{FD^{p}_{c}-IVA^{p}_{FD}-TESP^{p}}
\]

Con estas expresiones, el equilibrio base reproduce exactamente la SAM. El bien \emph{otros} queda excluido de ambos gravamenes, de modo que $t_{iva}^{p}=t_{iva,int}^{p}=\tau_{e}^{p}=0$ y $pq^{p}=pint^{p}=p^{p}$ para ese producto.

\subsection{Bloque de las firmas}

En este bloque, las firmas no se modelan como unidades físicas de producción ---función asignada al bloque de sectores---, sino como una institución que administra el flujo de renta del capital. Esta institución recibe ingresos, paga dividendos, impuestos y transferencias, y retiene el residuo como ahorro corporativo.

\textbf{1) Ingreso bruto de las firmas.} El ingreso inicial de las empresas se define como:

\[
ING^{F} = K^{F} p^{k} + CONT^{SOCF}, \qquad CONT^{SOCF}=\theta_{cs}^{F}\, CONT^{SOC}
\]

El primer componente corresponde a la renta del capital corporativo y el segundo, a la fracción de contribuciones sociales asignada a las firmas (identificada en la codificación de \texttt{gEcon} como \texttt{por\_cont\_soc\_f}).

\textbf{2) Transferencias a hogares (dividendos).} Las utilidades distribuidas se modelan como una proporción fija del ingreso bruto:

\[
TR^{EMP}=\theta_{tr}^{F}\, ING^{F}
\]

Esta identidad corresponde a \texttt{TREMP} en el archivo \texttt{.gcn}.

\textbf{3) Impuesto corporativo.} La recaudación sobre firmas se calcula sobre la base de renta del capital:

\[
FIRMTAX = fit\cdot\left(K^{F}p^{k}\right)
\]

donde $fit$ es la tasa del impuesto corporativo (primera categoría).

\textbf{4) Asignación institucional del capital.} La dotación de capital en manos de firmas corresponde a una participación fija del capital efectivamente utilizado en la economía:

\[
K^{F}=\kappa^{F}\, KD
\]

con $\kappa^{F}$ equivalente al parámetro exógeno \texttt{par\_k\_f} de calibración y $KD$ el capital utilizado definido en el bloque de equilibrio. En consecuencia, la base del impuesto corporativo varía con el grado de utilización del capital y no solo con su retorno.

\textbf{5) Ahorro corporativo residual.} El ahorro de firmas resulta como residuo luego de cumplir todos los pagos institucionales:

\[
SAV^{F}=ING^{F}-TR^{EMP}-PREST^{SOCF}-FIRMTAX-TR^{FG}
\]

\[
PREST^{SOCF}=\theta_{ps}^{F}\, ING^{F}, \qquad TR^{FG}=\theta_{fg}^{F}\, ING^{F}
\]

De manera equivalente, puede escribirse como la restricción presupuestaria de la institución de las firmas:

\[
ING^{F}=TR^{EMP}+PREST^{SOCF}+FIRMTAX+TR^{FG}+SAV^{F}
\]

En síntesis, este bloque transforma la renta del capital corporativo en distribución a hogares, pagos al gobierno y ahorro empresarial, cerrando el circuito financiero que luego alimenta el bloque ahorro--inversión.

\subsection{Bloque del gobierno}

El bloque de gobierno representa al sector público consolidado: recauda impuestos, recibe rentas y transferencias, ejecuta gasto público y determina el ahorro fiscal residual. No resuelve un problema de optimización explícito; opera como un sistema contable de cierre macroeconómico.

\textbf{1) Recaudación tributaria total.} La recaudación agregada se compone de seis fuentes:

\[
\begin{aligned}
TOTAL^{TAX}={}&\ VAT+PITAX+FIRMTAX \\
&+\sum_{s\in\mathrm{SEC}}\left(tps^{s}+tes^{s}\right)revenue^{s} \\
&+\sum_{p\in\mathrm{PRODS}}\left(TESP^{p}+ARANCEL^{p}\right)
\end{aligned}
\]

donde $VAT$ es el IVA agregado; $PITAX$, el impuesto a la renta de personas; $FIRMTAX$, el impuesto corporativo; y el término sectorial, los impuestos a la producción y otros gravámenes aplicados sobre las ventas brutas de cada actividad. En la aplicación para Chile, los impuestos a la producción incluyen tributos como timbres y estampillas, contribuciones territoriales, permisos de circulación, patentes y derechos municipales y no municipales, e impuestos a los juegos de azar, casinos e hipódromos. Por su parte, $ARANCEL^{p}$ corresponde a la recaudación aduanera y $TESP^{p}=\tau_{e}^{p}FD^{p}$, a los impuestos específicos sobre productos, que en la especificación vigente dejan de ser un monto exógeno de la SAM y pasan a responder al volumen de demanda final. Entre estos últimos se encuentran los impuestos de la Ley de Alcoholes y los impuestos específicos a los combustibles, el tabaco y las fuentes móviles. La recaudación agregada por este concepto suma el componente unitario por producto y el componente ad valorem $tes^{s}$ que la SAM registra a nivel de actividad. El IVA y los derechos de importación se mantienen desagregados en el modelo.

\textbf{2) IVA e impuesto a la renta de personas.} En forma desagregada:

\[
VAT=\sum_{p\in\mathrm{PRODS}} VAT^{p}
\]

\[
\begin{aligned}
PITAX={}&\ pit\cdot\Big(wL^{h}+p^{k}K^{h}+TR^{GOB}+TR^{EMP} \\
&\qquad\ +PS^{h}+FACTROW^{K}+FACTROW^{L}\Big)
\end{aligned}
\]

La primera identidad agrega el IVA recaudado por producto en el bloque \texttt{PRODUCTS}; la segunda aplica la tasa efectiva $pit$ sobre la base de ingresos del hogar.

\textbf{3) Ingreso total del fisco.} Los recursos del gobierno se escriben como:

\[
ING^{GOB}=p^{k}K^{g}+CONT^{SOCG}+TOTAL^{TAX}+TR^{FG}
\]

donde $p^{k}K^{g}$ corresponde a la renta del capital público, con $K^{g}=\kappa^{g}\,KD$; $CONT^{SOCG}$, a la fracción de contribuciones sociales que recibe el gobierno; y $TR^{FG}$, a las transferencias desde firmas.

\textbf{4) Transferencias del gobierno a hogares.} Se modelan como una proporción del ingreso fiscal:

\[
TR^{GOB}=\theta_{trgov}\,ING^{GOB}
\]

Esta regla replica el mecanismo de transferencias observado en la SAM.

\textbf{5) Gasto corriente y ahorro fiscal residual.} El ahorro del gobierno resulta de restar al ingreso total sus compromisos corrientes:

\[
SAV^{G}=ING^{GOB}-\left(\sum_{p\in\mathrm{PRODS}}pq^{p}GG^{p}+TR^{GOB}+PREST^{SOCG}\right)
\]

Equivalentemente,

\[
ING^{GOB}=\sum_{p\in\mathrm{PRODS}}pq^{p}GG^{p}+TR^{GOB}+PREST^{SOCG}+SAV^{G}
\]

El gasto del gobierno se valúra a precios de comprador $pq^{p}$, de modo que el sector público paga el IVA y los impuestos específicos asociados a sus compras. Como su volumen de compras permanece fijo en el nivel base, un alza de la tasa del IVA eleva simultáneamente la recaudación y el costo nominal del gasto público; el efecto neto sobre el ahorro fiscal recoge ambos movimientos.

En síntesis, el bloque fiscal concentra la recaudación proveniente de hogares, firmas, sectores y comercio exterior; financia gasto y transferencias; y su saldo $SAV^{G}$ interactúa con el cierre ahorro--inversión y con la necesidad de ahorro externo del modelo.

\subsection{Bloque del equilibrio}
\label{sec:equilibrio}

El bloque de equilibrio completa el cierre macroeconómico del modelo. Para ello, impone las restricciones de escasez y coherencia contable que deben cumplirse simultáneamente después de la optimización de los hogares, los sectores, los bloques CET--Armington, las firmas y el gobierno.

\textbf{1) Demanda agregada de trabajo.} La demanda total de trabajo se obtiene sumando el empleo utilizado por los conjuntos sectoriales definidos en la implementación y la fracción exógena asociada al resto del mundo:

\[
LD=\sum_{s\in \mathrm{SECTORS\_A}}L^{s}+\sum_{s\in \mathrm{SECTORS\_B}}L^{s}+\sum_{s\in \mathrm{SECTORS\_C}}L^{s}+L^{ROW}
\]

Esta identidad reemplaza el supuesto de pleno empleo: la cantidad efectivamente empleada surge de las decisiones de las firmas y no tiene por qué coincidir con la fuerza laboral total disponible.

\textbf{2) Participación laboral endógena.} La fuerza de trabajo de cada tipo responde al nivel de empleo, lo que recoge en forma reducida los efectos de trabajador desalentado y trabajador adicional:

\[
L^{FORCE}_{l}=L^{FORCE,data}_{l}\left(\frac{LD_{l}}{L^{data}_{l}}\right)^{\varepsilon}
\]

donde $L^{FORCE,data}_{l}$ y $L^{data}_{l}$ son la fuerza de trabajo y el empleo del escenario base, y $\varepsilon=0.42$. En la base, $LD_{l}=L^{data}_{l}$, de modo que la fuerza de trabajo coincide con el dato. Con este valor, una caída de 1\% del empleo reduce la fuerza de trabajo en 0.42\%, de modo que parte del ajuste se manifiesta como menor participación y no como mayor desempleo.

\textbf{3) Curva salarial.} El salario se determina mediante una curva salarial que vincula el salario con la tasa de ocupación de la fuerza de trabajo:

\[
w_{l}=w^{pot}_{l}\left(\frac{LD_{l}}{L^{FORCE}_{l}}\right)^{\phi}, \qquad w^{pot}_{l}=\left(1-u^{b}_{l}\right)^{-\phi}
\]

de modo que $w_{l}=1$ en el escenario base. Se adopta $\phi=2.0$. Como la curva está expresada en términos de la tasa de ocupación $(1-u)$ y no de la tasa de desempleo, su elasticidad respecto de $u$ es $-\phi\,u/(1-u)$, que con un desempleo agregado de 8.34\% equivale a $-0.18$. Este valor duplica aproximadamente la elasticidad central de $-0.1$ documentada por \cite{blanchflower_oswald_1994}, pero se encuentra dentro de la dispersión de esa literatura; su elección se discute en la subsección de calibración.

\textbf{4) Tasa de desempleo.} La diferencia entre fuerza laboral y empleo efectivo se resume en una identidad de desempleo:

\[
u_{l}=\frac{L^{FORCE}_{l}-LD_{l}}{L^{FORCE}_{l}}
\]

Con ello, los choques no se absorben únicamente mediante movimientos salariales, sino también mediante variaciones de la tasa de desempleo y de la participación. En la calibración base, las tasas de desempleo de los trabajos de baja, media y alta calificación son 7.6\%, 9.9\% y 7.8\%, respectivamente.

\textbf{5) Utilización endógena del capital.} El modelo no impone el pleno uso del stock de capital. El capital efectivamente utilizado, su grado de utilización y su retorno se determinan mediante:

\[
KD=\sum_{s\in \mathrm{SECTORS\_A}}K^{s}+\sum_{s\in \mathrm{SECTORS\_B}}K^{s}, \qquad
u_{k}=\frac{KD}{K^{TOT}}, \qquad
p^{k}=p^{k,pot}\,u_{k}^{\psi}
\]

con $\psi=0.2$. El stock $K^{TOT}$ deja de ser una restricción de recursos y pasa a representar la capacidad instalada: $u_{k}$ mide el grado en que esa capacidad se utiliza, y hacerlo más intensamente ---turnos adicionales, menor mantenimiento programado--- tiene un costo creciente en términos de retorno \cite{greenwood_hercowitz_huffman_1988}. La especificación es análoga a la curva salarial del mercado laboral, con $1/\psi$ como elasticidad de la oferta de servicios de capital. Bajo esta formulación, $p^{k}$ admite la interpretación de precio sombra del factor cuasi-fijo \cite{berndt_fuss_1986}.

La SAM registra la renta del capital efectivamente utilizado, de modo que $KD^{b}$ es observable pero $K^{TOT}$ no lo es. En versiones anteriores se imponía $u_{k}^{b}=1$, es decir, se identificaba el capital utilizado con la capacidad instalada. La versión vigente ancla la utilización base en $u_{k}^{b}=0.82$, valor tomado del uso de capacidad instalada que reporta el Indicador Mensual de Confianza Empresarial de ICARE y la Universidad Adolfo Ibáñez para diciembre de 2023 \cite{imce_2023}. La capacidad instalada y el retorno potencial se obtienen entonces de forma residual:

\[
K^{TOT}=\frac{KD^{b}}{u_{k}^{b}},
\qquad
p^{k,pot}=\left(u_{k}^{b}\right)^{-\psi}
\]

de modo que en el escenario base $p^{k}=1$ y el equilibrio reproduce la SAM. Conviene explicitar una propiedad de esta parametrización. Sustituyendo ambas expresiones en la curva de utilización se obtiene

\[
p^{k}=\left(u_{k}^{b}\right)^{-\psi}\left(\frac{KD}{K^{TOT}}\right)^{\psi}=\left(\frac{KD}{KD^{b}}\right)^{\psi},
\]

es decir, el retorno del capital depende únicamente de la variación del capital utilizado respecto de la base y no del nivel de $u_{k}^{b}$. Por lo tanto, el anclaje en 82\% es una renormalización: modifica el nivel reportado de utilización ---que pasa a interpretarse como porcentaje de la capacidad instalada y no como índice base 100--- sin alterar ninguna respuesta de equilibrio. Su valor es interpretativo y analítico: sitúa la economía a distancia de su techo de capacidad, hace comparables los resultados con las series de utilización empresariales \cite{corrado_mattey_1997} y es el paso previo necesario para imponer una restricción de capacidad $u_{k}\leq 1$ o una elasticidad $\psi$ creciente cerca del techo, extensión en la cual el nivel base sí alteraría las respuestas. Por construcción, en cualquier simulación se cumple $\Delta\ln p^{k}/\Delta\ln KD=\psi$, lo que sirve como verificación de la implementación.

Las dotaciones institucionales de capital ---hogares, firmas y gobierno--- escalan con $KD$. Como sus participaciones suman exactamente uno en la SAM (0.2158, 0.7643 y 0.0198, respectivamente) y el resto del mundo no recibe renta del capital, la renta asociada a la capacidad ociosa desaparece de los ingresos institucionales en la proporción correcta, sin generar ingresos residuales.

La motivación de este bloque es cuantitativa. Con capital plenamente utilizado y tecnología Cobb--Douglas en el valor agregado, la elasticidad del empleo respecto del nivel de actividad es exactamente $1/\beta_{L}$, independiente de la curva salarial. Con una participación del trabajo de $\beta_{L}=0.4597$, esa elasticidad es 2.18, lo que implica una respuesta del desempleo varias veces superior a la observada empíricamente.

\textbf{6) Numerario nominal.} Para cerrar el sistema de precios relativos se fija un numerario, mientras que el tipo de cambio se determina endógenamente. En esta versión del modelo, el tipo de cambio $exr$ es una variable de equilibrio que se ajusta para compatibilizar el balance externo con el nivel de ahorro externo fijado exógenamente.

En consecuencia, los precios relativos entre bienes domésticos e internacionales pueden variar. Esto permite que el sector externo absorba los choques macroeconómicos mediante cambios en la competitividad y en los términos de intercambio, en lugar de hacerlo mediante variaciones del ahorro externo.

\textbf{7) Participaciones fijas de los hogares.} La participación de los hogares en capital y prestaciones se calibra con parámetros exógenos, mientras que su dotación laboral depende del empleo efectivo:

\[
K^{h}=par_{k}^{h}\,KD, \qquad L^{h}=par_{l}^{h}\,LD, \qquad PS^{h}=psh^{data}
\]

Así, tanto el ingreso laboral como la renta del capital de los hogares se ajustan con el nivel de empleo y con el grado de utilización del capital.

\textbf{8) Coherencia comercial.} El sector externo queda cerrado dentro de los bloques CET y Armington: las ventas externas $EX^{p}$, las ventas domésticas $DS^{p}$, las importaciones $M^{p}$ y el bien compuesto $Q^{p}$ se determinan endógenamente a partir de precios relativos, sin necesidad de introducir rentas residuales adicionales en este bloque.

\subsection{Calibración de los mercados de factores}

Los parámetros $\phi$, $\psi$ y $\varepsilon$ se eligen conjuntamente para que la respuesta del desempleo al producto sea coherente con la evidencia empírica. Log-linealizando las demandas de factores de la tecnología Cobb--Douglas con insumos Leontief, junto con la curva salarial, la curva de participación y la curva de utilización del capital, se obtiene:

\[
\frac{\hat{l}}{\hat{v}}=\frac{1+\psi}{1+\psi\beta_{L}+\phi(1-\varepsilon)\beta_{K}},
\qquad
\frac{\Delta u}{\hat{v}}=(1-u)\,(1-\varepsilon)\,\frac{\hat{l}}{\hat{v}}
\]

donde $\hat{l}$ y $\hat{v}$ son las variaciones porcentuales del empleo y de la actividad, y $\beta_{L}=0.4597$ y $\beta_{K}=0.5403$ son las participaciones agregadas del trabajo y del capital en el valor agregado de la SAM. La segunda expresión corresponde al coeficiente de Okun implícito del modelo \cite{okun_1962} y puede reescribirse como $(1-u)(\eta_{L}-\eta_{LF})$, donde $\eta_{L}$ y $\eta_{LF}$ son las elasticidades del empleo y de la fuerza de trabajo respecto del producto.

Dos propiedades de estas expresiones orientan la calibración. Primero, cuando el capital se mantiene plenamente utilizado ($\psi\to\infty$), $\hat{l}/\hat{v}\to 1/\beta_{L}=2.18$ con independencia de $\phi$: ningún valor de la curva salarial corrige la sobrerreacción del desempleo. Segundo, aun con oferta de capital perfectamente elástica ($\psi=0$) y fuerza de trabajo fija ($\varepsilon=0$), el coeficiente de Okun no puede ser inferior a $(1-u)/(1+\phi\beta_{K})$, que para valores razonables de $\phi$ excede la evidencia. El margen de participación es, por tanto, necesario.

Como referencia empírica se utiliza el informe de \cite{clapes_okun}, que estima la ley de Okun en diferencias para Chile con datos trimestrales de 2010 a 2023. La especificación estática arroja un coeficiente de $-0.281$. La especificación dinámica arroja un coeficiente de impacto de $-0.207$ y una persistencia del desempleo de 0.572, lo que implica un multiplicador de largo plazo de $-0.367$. Dado que un modelo estático compara equilibrios con ajuste completo, se adopta como objetivo la banda 0.28--0.37 y, dentro de ella, preferentemente su extremo superior.

La Tabla~\ref{tab:calib-factores} resume los valores adoptados. Con ellos, la aproximación analítica entrega un coeficiente de Okun de 0.371, una elasticidad empleo--producto de 0.70 y una elasticidad fuerza de trabajo--producto de 0.29. En las simulaciones, el coeficiente medido respecto del valor agregado real se ubica en 0.326--0.327 para los tres choques de base amplia (IVA, PIT e IDPC) y desciende a 0.276 en el alza de los impuestos específicos. La diferencia con la aproximación analítica refleja la composición sectorial de cada choque, que la expresión agregada no captura: cuando la medida se concentra en unos pocos productos, la intensidad laboral de los sectores afectados determina la respuesta agregada del empleo.

\begin{table}[H]
\centering
\caption{Parámetros de los mercados de factores}
\label{tab:calib-factores}
\small
\begin{tabular}{|l|c|c|p{6.5cm}|}
  \hline
  Parámetro & Símbolo & Valor & Interpretación \\
  \hline
  Curva salarial & $\phi$ & 2.0 & Elasticidad salario--desempleo implícita de $-0.18$ \\
  \hline
  Utilización del capital & $\psi$ & 0.2 & Elasticidad de la oferta de servicios de capital de 5 \\
  \hline
  Participación laboral & $\varepsilon$ & 0.42 & Elasticidad de la fuerza de trabajo respecto del empleo \\
  \hline
  Utilización base del capital & $u_{k}^{b}$ & 0.82 & Uso de capacidad instalada, IMCE ICARE--UAI, diciembre 2023. Renormaliza el nivel reportado sin alterar las respuestas de equilibrio \\
  \hline
\end{tabular}
\end{table}

Esta calibración tiene limitaciones que conviene explicitar. El informe utilizado como referencia no ha sido sometido a revisión de pares; su regresión en diferencias interanuales sobre datos trimestrales induce autocorrelación en los errores, y su muestra incluye los trimestres de la pandemia, que pueden tener una influencia considerable sobre el coeficiente estimado. El parámetro $\psi$ carece de un ancla empírica directa para Chile, donde no existe una serie oficial de utilización de capacidad por sector, y $\varepsilon$ se obtiene de la identidad de Okun y no de una estimación propia de $\eta_{LF}$. El nivel base $u_{k}^{b}$ proviene de una encuesta de opinión empresarial de cobertura principalmente industrial, por lo que su extensión al conjunto de la economía es una aproximación; su efecto, sin embargo, se limita a la interpretación de los niveles reportados. Por estas razones, los tres parámetros se incluyen en el análisis de sensibilidad. Finalmente, la curva salarial y la curva de participación son dispositivos de forma reducida que replican la elasticidad observada, no una teoría estructural del desempleo.


\subsection{Bloque de ahorro--inversión y cierre externo}

El cierre ahorro--inversión no es neoclásico. La inversión del gobierno permanece fija en su nivel base, mientras que la inversión de los hogares y las firmas varía endógenamente en función de su ahorro, con una respuesta amortiguada. En consecuencia, el ajuste macroeconómico combina cambios en la inversión privada y en los precios relativos, especialmente a través del tipo de cambio, mientras el financiamiento externo permanece constante.

\textbf{1) Inversión pública fija e inversión privada variable.} Primero se define el costo nominal de la canasta base de inversión:

\[
NOM\_I^{BASE}=\sum_{p\in\mathrm{PRODS}} I^{data,p}\,pq^{p}
\]

Luego, el componente gubernamental se mantiene fijo mediante:

\[
inv\_adj^{G}=1
\]

En cambio, los componentes de hogares y firmas se ajustan en proporción a su ahorro, con un parámetro de amortiguación $\eta$:

\[
inv\_adj^{H}=\left(\frac{SAV}{sh\_i^{H}\,NOM\_I^{BASE}}\right)^{\eta},
\qquad
inv\_adj^{F}=\left(\frac{SAV^{F}}{sh\_i^{F}\,NOM\_I^{BASE}}\right)^{\eta}
\]

En la calibración utilizada se fija $\eta=0.5$\footnote{Con este valor, la inversión de hogares y firmas responde a cambios en el ahorro con una elasticidad menor que uno y, por tanto, de manera amortiguada.}.

El multiplicador agregado de inversión que entra al bloque de productos es un promedio ponderado de esos tres componentes:

\[
inv\_adj=sh\_i^{H}inv\_adj^{H}+sh\_i^{F}inv\_adj^{F}+sh\_i^{G}inv\_adj^{G}
\]

Así, la demanda de inversión por producto no permanece completamente fija. La inversión pública queda anclada a la canasta calibrada, mientras que la inversión de los hogares y las firmas puede expandirse o contraerse según su ahorro disponible, con una elasticidad acotada por $\eta$.

\textbf{2) Requerimiento total de financiamiento.} El financiamiento total asociado a la inversión agregada y a la variación de existencias viene dado por:

\[
TOTAL^{INV}=NOM\_I^{BASE}\,inv\_adj+\sum_{p\in\mathrm{PRODS}}VEXIST^{p}p^{p}
\]

De este modo, el costo nominal de la inversión total depende del ajuste conjunto entre inversión privada variable, inversión pública fija y variación de existencias.

\textbf{3) Cierre externo con ahorro externo fijo.} En la versión actual del modelo, el ahorro externo se considera exógeno y se fija en su nivel de referencia:

\[
SAV^{EXT}=SAV_{\text{target}}^{EXT}
\]

Dado este cierre, el equilibrio macroeconómico se logra mediante ajustes en el tipo de cambio real, el cual se determina endógenamente para compatibilizar el balance entre ahorro e inversión agregados.

En este contexto, los choques fiscales o distributivos no pueden absorberse mediante variaciones del financiamiento externo. Por ello, el ajuste ocurre a través de cambios en los precios relativos, que afectan la competitividad externa, las decisiones de consumo e inversión y la composición del gasto entre bienes domésticos e importados.

Este cierre implica que la restricción externa es activa y que la economía no puede expandir su nivel de absorción más allá de lo permitido por su capacidad de financiamiento externo.

\section{Metodología de evaluación del impacto fiscal}

La evaluación fiscal compara el equilibrio base con el obtenido después de aplicar cada medida tributaria. Además de calcular la variación total de la recaudación, el procedimiento distingue los efectos asociados a la tasa y a la base real del efecto originado por la revaloración de los precios.

\subsection{Métrica de impacto recaudatorio}

Para cada instrumento tributario \(i\), la recaudación en los escenarios base y simulado se expresa como porcentaje del PIB nominal del escenario base:

\[
r_i^k=100\frac{R_i^k}{\mathrm{PIB}_{nom}^{b}},
\qquad k\in\{b,s\},
\]

donde \(b\) identifica el equilibrio base y \(s\) el equilibrio simulado. Por tanto, el impacto fiscal se mide en puntos porcentuales del PIB nominal base:

\[
\Delta r_i=100\frac{R_i^s-R_i^b}{\mathrm{PIB}_{nom}^{b}}.
\]

El uso de un denominador fijo evita que el cambio reportado combine variaciones de recaudación con movimientos simultáneos del PIB. Asimismo, se utiliza el PIB nominal porque tanto la recaudación como sus bases valoradas se expresan en unidades corrientes del numerario del modelo.

\subsection{Descomposición entre tasa, base real y precio}
\label{sec:metodologia-descomp}

La recaudación de cada instrumento se representa como la suma de productos entre una tasa \(\tau\), un precio \(P\) y una cantidad o base real \(Q\):

\[
R_i=\sum_j \tau_{ij}P_{ij}Q_{ij},
\]

donde \(j\) identifica productos, sectores o componentes de ingreso, según el impuesto analizado. Se aplica la descomposición de Shapley para un producto de tres factores \cite{shapley_1953}. Definiendo las variaciones entre escenarios como \(\Delta x=x^s-x^b\) y los promedios como \(\bar{x}=(x^s+x^b)/2\), los efectos se calculan mediante:

\begin{align*}
E_i^{tasa} &= \sum_j \Delta\tau_{ij}\left(\bar P_{ij}\bar Q_{ij}+\tfrac{1}{12}\Delta P_{ij}\Delta Q_{ij}\right),\\
E_i^{base} &= \sum_j \Delta Q_{ij}\left(\bar\tau_{ij}\bar P_{ij}+\tfrac{1}{12}\Delta\tau_{ij}\Delta P_{ij}\right),\\
E_i^{precio} &= \sum_j \Delta P_{ij}\left(\bar\tau_{ij}\bar Q_{ij}+\tfrac{1}{12}\Delta\tau_{ij}\Delta Q_{ij}\right).
\end{align*}

La fórmula simétrica de punto medio, sin los términos de orden superior, omite exactamente \(\tfrac{1}{4}\Delta\tau\,\Delta P\,\Delta Q\); los términos \(\tfrac{1}{12}\) reparten esa interacción en partes iguales entre los tres efectos \cite{sun_1998}. En consecuencia, la descomposición es exacta:

\[
\Delta R_i=E_i^{tasa}+E_i^{base}+E_i^{precio},
\]

y la discrepancia numérica entre ambos lados es del orden de la precisión de máquina (inferior a \(10^{-14}\) puntos del PIB en todas las simulaciones), por lo que no se reporta un residuo. El reporte presenta también el efecto conjunto de tasa y base:

\[
E_i^{tasa+base}=E_i^{tasa}+E_i^{base}.
\]

Esta agrupación permite separar el efecto precio del efecto conjunto de la medida tributaria y de la respuesta real de la base. Todos los componentes se dividen por el PIB nominal base y se expresan en puntos porcentuales del PIB.

Para los impuestos específicos sobre productos, la recaudación es un producto de solo dos factores, $R=\sum_{p}\tau_{e}^{p}FD^{p}$, por lo que la descomposición se reduce a su versión bilineal exacta:

\[
E^{tasa}=\sum_{p}\Delta\tau_{e}^{p}\,\overline{FD}^{p},
\qquad
E^{base}=\sum_{p}\Delta FD^{p}\,\bar{\tau}_{e}^{p},
\qquad
E^{precio}\equiv 0.
\]

La anulación del efecto precio no es una simplificación del método, sino una propiedad del instrumento: un impuesto unitario no se revaloriza con los precios nominales. Bajo la representación ad valorem previa, la recaudación específica heredaba un efecto precio que no corresponde al diseño legal del tributo.

\subsection{Aplicación por instrumento tributario}

La identificación de tasa, precio y base real depende de la estructura de cada impuesto. La Tabla~\ref{tab:descomposicion-fiscal} resume la correspondencia utilizada en la implementación.

\begin{table}[H]
\centering
\caption{Componentes utilizados en la descomposición del impacto fiscal}
\label{tab:descomposicion-fiscal}
\small
\begin{tabular}{|p{2.2cm}|p{2.2cm}|p{3.2cm}|p{6.1cm}|}
  \hline
  Instrumento & Tasa & Precio & Base real o cantidad \\
  \hline
  IVA & \(t_{iva}^p\) sobre demanda final y \(t_{iva,int}^p\) sobre uso intermedio & Precio con específico \(p^p+\tau_e^p\) para la demanda final y precio básico \(p^p\) para el uso intermedio & Demanda final física \(FD^p\) y demanda intermedia física \(XI^p\), disponibles directamente como variables del modelo. \\
  \hline
  Impuestos específicos sobre productos & \(\tau_e^p\), monto por unidad física & No aplica: la recaudación no depende de precios & Demanda final física \(FD^p\). La descomposición se reduce a dos factores y el efecto precio es idénticamente nulo. \\
  \hline
  IDPC & \(fit\) & Retorno del capital \(p^k\) & Capital corporativo \(K^F=\kappa^F KD\), que varía con la utilización del capital. Dado que \(p^k=p^{k,pot}u_k^{\psi}\), los efectos precio y base del IDPC guardan la razón \(\psi\). \\
  \hline
  PIT & \(pit\) & \(p^k\) para el capital, \(w^l\) para el trabajo y un valor unitario para las transferencias & Capital de los hogares \(K^h\), empleo por tipo \(L^{h,l}\) y transferencias recibidas. Los efectos de las subbases se calculan por separado y luego se agregan. \\
  \hline
  Impuestos a la producción y gravámenes sectoriales & \(tps^s\) y \(tes^s\) & Índice de precios de producción construido a partir de \(pz^p\) & Producción sectorial implícita, obtenida al deflactar los ingresos sectoriales. \\
  \hline
  Aranceles & \(tar\_rate^p\) & Tipo de cambio \(exr\) & Importaciones \(M^p\) valoradas a precios mundiales. Si la tasa no cambia, el efecto tasa es cero. \\
  \hline
\end{tabular}
\end{table}

\subsection{Interpretación de los efectos}

El efecto tasa aproxima el rendimiento mecánico del cambio tributario, manteniendo constantes la valoración y la base real. El efecto base recoge la respuesta de las cantidades inducida por el nuevo equilibrio ---por ejemplo, cambios en el consumo, la producción, el empleo o las importaciones---. El efecto precio refleja la revaloración nominal de la base provocada por movimientos de los precios relativos, los salarios, el retorno del capital o el tipo de cambio.

Esta separación es relevante porque el efecto precio depende del numerario y del cierre macroeconómico. En consecuencia, no debe interpretarse de manera aislada como una proyección del flujo de caja fiscal. El efecto base contiene la respuesta real de los agentes, mientras que el efecto tasa constituye el punto de comparación con una estimación estática.

\subsection{Medición de la actividad real}

El PIB real se calcula por el lado del gasto, valorando los componentes de demanda final a precios del escenario base. En versiones anteriores del modelo, el IVA se representaba en el equilibrio de cada mercado de productos como unidades de cantidad proporcionales a la absorción gravada. Por ello, cuando cambiaba la tasa del IVA, la variación de esas unidades se registraba en el PIB real como un cambio de volumen y no de precio, lo que sesgaba la medida en simulaciones que modificaban el IVA \cite{sna_2008}. La reespecificación descrita en la Sección~\ref{sec:productos} elimina ese sesgo por construcción: como $Q^{p}=TD\_Dom^{p}$ en unidades físicas y los impuestos operan sobre precios, un alza de la tasa no genera volumen ficticio. La verificación numérica es concluyente: en el choque del IVA la base real implícita de los impuestos indirectos pasa de aumentar 2.74\% a caer 1.21\%, acompañando la contracción de los volúmenes, y la brecha entre ambas medidas de actividad se reduce de 0.47 a 0.035 puntos porcentuales.

Por este motivo se reporta también el valor agregado real a costo de factores:

\[
VA^{real}=\sum_{s}\frac{va^{b}_{s}}{V^{b}_{s}}\,V^{s}_{s},
\]

donde $va^{b}_{s}$ es el valor agregado nominal del sector en la base y $V_{s}$ su nivel de actividad. La expresión es válida porque los insumos intermedios son proporcionales a la actividad, de modo que el valor agregado por unidad de producto es constante a precios base. Bajo la especificación vigente, ambas medidas difieren en 0.035 puntos porcentuales en el choque del IVA, 0.009 en el del PIT, 0.004 en el del IDPC y 0.012 en el de los impuestos específicos. Con la especificación anterior, la brecha del choque del IVA era de 0.47 puntos, un orden de magnitud mayor. El valor agregado real se mantiene como medida de referencia porque es invariante a la valoración de la demanda final y permite comparar los tres choques sobre una misma base. El coeficiente de Okun se evalúa respecto del valor agregado real.


\section{Simulaciones de políticas tributarias}
\label{sec:simulaciones}

Esta sección evalúa el impacto macroeconómico y fiscal de cuatro políticas tributarias: el incremento del IVA, el aumento del impuesto a la renta de personas (PIT), la disminución del impuesto a las firmas y el alza de los impuestos específicos sobre productos. Para cada una se reportan el PIB real medido por el gasto y el valor agregado real a costo de factores, la utilización de los factores y la descomposición del impacto fiscal.

La reespecificación de los impuestos indirectos tiene un efecto visible en estos resultados. Bajo el tratamiento anterior, la recaudación de impuestos específicos era estrictamente constante ante cualquier choque, de modo que la fila correspondiente registraba ceros en las tres descomposiciones. Con el impuesto unitario sobre demanda final, esa recaudación responde a los volúmenes: aporta -0.0139 puntos del PIB en el alza del IVA, -0.0056 en la del PIT y +0.0021 en la rebaja del IDPC, siempre con efecto precio exactamente nulo. Además, el alza de los propios impuestos específicos pasa a ser un choque simulable, lo que se desarrolla en la Sección~\ref{sec:sim-tesp}.

\subsection{Incremento del IVA}
\label{sec:sim-iva}

En esta sección se evalúa el impacto macroeconómico de un aumento de un punto porcentual en la tasa del IVA, equivalente a pasar de una tasa de referencia de 19\% a 20\%\footnote{En la implementación vigente, las tasas efectivas por producto $t_{iva}^{p}$ ---expresadas sobre el precio sin impuesto--- se multiplican por el factor $20/19$ definido en el archivo de configuración. Como la cuña es multiplicativa, esto equivale a elevar la tasa legal de 19\% a 20\% preservando la estructura relativa de tasas efectivas calibrada desde la SAM.}. Por tanto, el choque corresponde a un incremento proporcional de aproximadamente 5.26\% en las tasas efectivas calibradas.

La Figura~\ref{fig:vat-transmission} resume los principales canales de transmisión de esta política. Bajo el cierre con tipo de cambio flexible y ahorro externo fijo, el choque aumenta inicialmente el precio al consumidor, reduce el ingreso real de los hogares y contrae el consumo. La caída de la absorción interna se transmite luego a la producción, al empleo y a la utilización del capital. Al mismo tiempo, la menor demanda de importaciones reduce la presión externa y desencadena un ajuste cambiario, mientras que el aumento de la tasa eleva directamente la recaudación por IVA.

\begin{figure}[H]
\centering
\includegraphics[width=\textwidth]{diagrams/diagram_02_aumento_iva.png}
\caption{Diagrama conceptual de transmisión para el aumento del IVA.}
\label{fig:vat-transmission}
\end{figure}

\begin{table}[H]
\centering
\caption{Resultados agregados del incremento del IVA}
\begin{tabular}{|l|r|r|r|}
  \hline
   & Escenario inicial & Escenario simulado & Cambio \\
  \hline
  PIB nominal & 282097.80 & 279515.35 & -0.9154\% \\
  \hline
  PIB real (gasto) & 282097.80 & 279430.61 & -0.9455\% \\
  \hline
  Valor agregado real (costo de factores) & 249320.66 & 247050.60 & -0.9105\% \\
  \hline
\end{tabular}
\end{table}

\begin{table}[H]
\centering
\caption{Componentes de demanda del incremento del IVA}
\begin{tabular}{|l|r|r|}
  \hline
   & Cambio real & Cambio nominal \\
  \hline
  Consumo & -1.2603\% & -1.1367\% \\
  \hline
  Inversión & -0.6549\% & -0.5657\% \\
  \hline
  Gasto de gobierno & 0.0000\% & -- \\
  \hline
  Exportaciones & -1.3389\% & -1.6619\% \\
  \hline
  Importaciones & -1.2997\% & -1.6228\% \\
  \hline
\end{tabular}
\end{table}

\begin{table}[H]
\centering
\caption{Comercio exterior del incremento del IVA}
\begin{tabular}{|l|r|r|r|}
  \hline
   & Escenario inicial & Escenario simulado & Cambio \\
  \hline
  Exportaciones & 86815.42 & 85372.66 & -1442.75 \\
  \hline
  Importaciones & 84351.36 & 82982.49 & -1368.88 \\
  \hline
  Saldo comercial & 2464.05 & 2390.18 & -73.88 \\
  \hline
\end{tabular}
\end{table}

\begin{table}[H]
\centering
\caption{Mercado laboral del incremento del IVA}
\begin{tabular}{|l|r|r|r|}
  \hline
   & Escenario inicial & Escenario simulado & Cambio \\
  \hline
  Desempleo, baja calificación & 7.60\% & 7.93\% & 0.3313 p.p. \\
  \hline
  Desempleo, calificación media & 9.90\% & 10.21\% & 0.3145 p.p. \\
  \hline
  Desempleo, alta calificación & 7.80\% & 8.09\% & 0.2879 p.p. \\
  \hline
\end{tabular}
\end{table}

\begin{table}[H]
\centering
\caption{Utilización de factores y respuesta agregada del empleo del incremento del IVA}
\small
\begin{tabular}{|l|r|}
  \hline
  Indicador & Valor \\
  \hline
  Utilización del capital, escenario simulado (\% de capacidad instalada; base 82.00\%) & 81.0076\% \\
  \hline
  Capital utilizado ($KD$) & -1.2102\% \\
  \hline
  Retorno del capital ($p^{k}$) & -0.2432\% \\
  \hline
  Desempleo agregado (base 8.3359\%) & 8.6327\% \\
  \hline
  Elasticidad empleo--producto $\eta_{L}$ (respecto del PIB real) & 0.5900 \\
  \hline
  Elasticidad fuerza de trabajo--producto $\eta_{LF}$ (respecto del PIB real) & 0.2483 \\
  \hline
  Coeficiente de Okun respecto del PIB real & 0.3139 \\
  \hline
  Coeficiente de Okun respecto del valor agregado real & 0.3260 \\
  \hline
\end{tabular}
\end{table}

\begin{table}[H]
\centering
\caption{Variación de recaudación tributaria del incremento del IVA (porcentaje del PIB nominal inicial)}
\small
\begin{tabular}{|l|r|r|r|}
  \hline
   & Base (\% del PIB) & Simulado (\% del PIB) & Cambio (p.p. del PIB) \\
  \hline
  Recaudación total & 18.1916 & 18.3671 & 0.1755 \\
  \hline
  IDPC (\textit{FIRMTAX}) & 4.7622 & 4.6932 & -0.0691 \\
  \hline
  PIT & 1.8103 & 1.7897 & -0.0206 \\
  \hline
  IVA & 8.6254 & 8.9320 & 0.3067 \\
  \hline
  Impuestos a la producción & 1.4745 & 1.4560 & -0.0185 \\
  \hline
  Impuestos específicos & 0.9727 & 0.9588 & -0.0139 \\
  \hline
  Aranceles & 0.5465 & 0.5373 & -0.0091 \\
  \hline
\end{tabular}
\end{table}

\begin{table}[H]
\centering
\caption{Descomposición del impacto fiscal del incremento del IVA (puntos del PIB nominal inicial)}
\label{tab:iva-descomposicion}
\small
\resizebox{\textwidth}{!}{%
\begin{tabular}{|l|r|r|r|r|r|}
  \hline
  Impuesto & Cambio observado & Efecto tasa & Efecto base & Tasa más base & Efecto precio \\
  \hline
  IVA & 0.3067 & 0.4503 & -0.1088 & 0.3414 & -0.0348 \\
  \hline
  IDPC (\textit{FIRMTAX}) & -0.0691 & 0.0000 & -0.0576 & -0.0576 & -0.0115 \\
  \hline
  PIT & -0.0206 & 0.0000 & -0.0133 & -0.0133 & -0.0073 \\
  \hline
  Impuestos a la producción & -0.0185 & 0.0000 & -0.0123 & -0.0123 & -0.0063 \\
  \hline
  Impuestos específicos & -0.0139 & 0.0000 & -0.0139 & -0.0139 & 0.0000 \\
  \hline
  Aranceles & -0.0091 & 0.0000 & -0.0073 & -0.0073 & -0.0018 \\
  \hline
  Total & 0.1755 & 0.4503 & -0.2132 & 0.2371 & -0.0616 \\
  \hline
\end{tabular}%
}
\end{table}

El aumento de un punto porcentual en la tasa del IVA contrae la actividad agregada: el valor agregado real a costo de factores disminuye 0.9105\% y el PIB real medido por el gasto, 0.9455\%. Ambas medidas difieren ahora en 0.035 puntos porcentuales, frente a los 0.47 puntos que se obtenían bajo la especificación anterior. La base real implícita de los impuestos indirectos cae 1.2116\%, en lugar de aumentar 2.74\% como ocurría cuando el impuesto absorbía unidades físicas del bien compuesto. Este es el diagnóstico más directo de que el sesgo de medición identificado en la Sección~\ref{sec:productos} quedó corregido: el alza de la tasa ya no se registra como volumen.

El PIB nominal cae 0.9154\%, es decir, menos que el PIB real. La relación se invierte respecto de la especificación anterior y es la esperada: el PIB a precios de mercado incorpora los impuestos sobre productos, de modo que una cuña mayor sostiene el valor nominal mientras los volúmenes se contraen. El mismo patrón aparece en el consumo, con una caída real de 1.2603\% y nominal de 1.1367\%.

La contracción de la actividad es también sustancialmente menor que la estimada antes: 0.91\% del valor agregado real frente a 1.32\%. La diferencia obedece al crédito pleno sobre el uso intermedio. En la especificación anterior, el impuesto encarecía implícitamente los insumos y distorsionaba las decisiones de producción; con la cadena de débitos y créditos operativa, el IVA recae sobre la demanda final y preserva la eficiencia productiva, en el sentido de \cite{diamond_mirrlees_1971}. El resultado no es un artificio de calibración: es la propiedad que justifica la arquitectura misma del impuesto al valor agregado \cite{ebrill_etal_2001}.

Por componentes, el consumo real cae 1.2603\%, la inversión real disminuye 0.6549\% y el gasto real del gobierno permanece constante. En el frente externo, las exportaciones reales caen 1.3389\% y las importaciones reales, 1.2997\%, por lo que el saldo comercial disminuye de 2464.05 a 2390.18. Esta combinación es coherente con el mecanismo ilustrado en la figura: la menor absorción interna reduce las importaciones, pero el ajuste de los precios relativos debilita en mayor medida las exportaciones. El consumo es ahora el componente más afectado, por encima de las exportaciones, lo que refleja que la cuña recae sobre la demanda final doméstica y no sobre los costos de producción.

En los mercados de factores, la tasa de desempleo aumenta 0.3313 puntos porcentuales para el trabajo de baja calificación, 0.3145 para el de calificación media y 0.2879 para el de alta calificación. La utilización del capital cae de 82.00\% a 81.01\% de la capacidad instalada, con una contracción de 1.21\% del capital utilizado, y el retorno del capital disminuye 0.24\%. El coeficiente de Okun respecto del valor agregado real es 0.326, prácticamente idéntico al de los choques de PIT e IDPC. Bajo la especificación anterior, el coeficiente medido respecto del PIB real ascendía a 0.503 por el sesgo del denominador; corregida la medición, ambos indicadores convergen a 0.314 y 0.326, lo que elimina la necesidad de advertir sobre esa discrepancia.

La variación nominal observada de la recaudación total es de 0.1755 puntos del PIB inicial. Al excluir el efecto precio, el efecto conjunto de tasa y base alcanza 0.2371 puntos del PIB. Este monto se compone de un efecto tasa mecánico de 0.4503 puntos, asociado íntegramente al IVA, y de un efecto base de -0.2132 puntos, que absorbe el 47\% del rendimiento mecánico. La contracción reduce la base del propio IVA en 0.1088 puntos y la del IDPC en 0.0576 puntos, esta última por la menor utilización del capital corporativo. Las bases del PIT, los impuestos a la producción, los impuestos específicos y los aranceles disminuyen en 0.0133, 0.0123, 0.0139 y 0.0073 puntos, respectivamente. La pérdida en impuestos específicos ---nula por construcción en la especificación anterior--- es enteramente de cantidad y refleja el menor consumo físico de los bienes gravados.

El efecto precio agregado es de -0.0616 puntos del PIB, y sus principales componentes corresponden al IVA y al IDPC, con -0.0348 y -0.0115 puntos, respectivamente. Su magnitud se reduce a dos tercios de la estimada bajo la especificación anterior, de modo que el resultado recaudatorio depende menos de la revaloración nominal de las bases. Aún así, dado que este componente depende del numerario y del ajuste de los precios relativos, la evaluación fiscal de la política se concentra en la suma de los efectos tasa y base.

\textbf{Conclusión del incremento del IVA:} el aumento de la tasa de referencia de 19\% a 20\% genera un efecto conjunto de tasa y base equivalente a 0.2371 puntos del PIB inicial. El rendimiento mecánico de la mayor tasa es atenuado en 47\% por la contracción de la base del IVA y de las demás bases tributarias, en particular la del impuesto corporativo. Este aumento de la recaudación tiene como contrapartida una caída de 0.91\% del valor agregado real, de la utilización del capital y del empleo, junto con un deterioro del saldo comercial. Relativo al resto de las medidas evaluadas, el IVA resulta el instrumento menos costoso por punto recaudado: 3.8 puntos de valor agregado real por punto del PIB de recaudación de tasa más base, frente a 5.2 en los impuestos específicos y 5.4 en el impuesto a la renta de personas. Esta jerarquía se invierte respecto de la versión anterior del modelo, en la que el IVA apareció como el más costoso (6.3), y la explicación es la desgravación del uso intermedio.

\subsection{Aumento del impuesto a la renta de personas (PIT)}

Esta simulación evalúa un aumento de 10\% en la tasa efectiva del impuesto a la renta de personas, representada en el modelo por el parámetro \textit{pit}. Bajo el cierre con tipo de cambio flexible y ahorro externo fijo, el choque reduce el ingreso disponible de los hogares y contrae el consumo. Al disminuir el ahorro de los hogares, también debilita parcialmente la inversión privada. La Figura~\ref{fig:pit-transmission} resume este mecanismo: el menor gasto privado comprime la absorción interna, reduce la producción y la utilización de los factores, y activa un ajuste externo mediante los precios relativos.

\begin{figure}[H]
\centering
\includegraphics[width=\textwidth]{diagrams/diagram_03_aumento_pit.png}
\caption{Diagrama conceptual de transmisión para el aumento del impuesto a la renta de personas (PIT).}
\label{fig:pit-transmission}
\end{figure}

\begin{table}[H]
\centering
\caption{Resultados agregados del aumento del PIT}
\begin{tabular}{|l|r|r|r|}
  \hline
   & Escenario inicial & Escenario simulado & Cambio \\
  \hline
  PIB nominal & 282097.80 & 280195.92 & -0.6742\% \\
  \hline
  PIB real (gasto) & 282097.80 & 280787.82 & -0.4644\% \\
  \hline
  Valor agregado real (costo de factores) & 249320.66 & 248184.01 & -0.4559\% \\
  \hline
\end{tabular}
\end{table}

\begin{table}[H]
\centering
\caption{Componentes de demanda del aumento del PIT}
\begin{tabular}{|l|r|r|}
  \hline
   & Cambio real & Cambio nominal \\
  \hline
  Consumo & -0.6335\% & -0.8359\% \\
  \hline
  Inversión & -0.2541\% & -0.4497\% \\
  \hline
  Gasto de gobierno & 0.0000\% & -- \\
  \hline
  Exportaciones & -0.6146\% & -0.7846\% \\
  \hline
  Importaciones & -0.5683\% & -0.7384\% \\
  \hline
\end{tabular}
\end{table}

\begin{table}[H]
\centering
\caption{Comercio exterior del aumento del PIT}
\begin{tabular}{|l|r|r|r|}
  \hline
   & Escenario inicial & Escenario simulado & Cambio \\
  \hline
  Exportaciones & 86815.42 & 86134.25 & -681.17 \\
  \hline
  Importaciones & 84351.36 & 83728.48 & -622.89 \\
  \hline
  Saldo comercial & 2464.05 & 2405.77 & -58.28 \\
  \hline
\end{tabular}
\end{table}

\begin{table}[H]
\centering
\caption{Mercado laboral del aumento del PIT}
\begin{tabular}{|l|r|r|r|}
  \hline
   & Escenario inicial & Escenario simulado & Cambio \\
  \hline
  Desempleo, baja calificación & 7.60\% & 7.76\% & 0.1577 p.p. \\
  \hline
  Desempleo, calificación media & 9.90\% & 10.05\% & 0.1492 p.p. \\
  \hline
  Desempleo, alta calificación & 7.80\% & 7.95\% & 0.1479 p.p. \\
  \hline
\end{tabular}
\end{table}

\begin{table}[H]
\centering
\caption{Utilización de factores y respuesta agregada del empleo del aumento del PIT}
\small
\begin{tabular}{|l|r|}
  \hline
  Indicador & Valor \\
  \hline
  Utilización del capital, escenario simulado (\% de capacidad instalada; base 82.00\%) & 81.5033\% \\
  \hline
  Capital utilizado ($KD$) & -0.6057\% \\
  \hline
  Retorno del capital ($p^{k}$) & -0.1214\% \\
  \hline
  Desempleo agregado (base 8.3359\%) & 8.4846\% \\
  \hline
  Elasticidad empleo--producto $\eta_{L}$ (respecto del PIB real) & 0.6022 \\
  \hline
  Elasticidad fuerza de trabajo--producto $\eta_{LF}$ (respecto del PIB real) & 0.2531 \\
  \hline
  Coeficiente de Okun respecto del PIB real & 0.3203 \\
  \hline
  Coeficiente de Okun respecto del valor agregado real & 0.3263 \\
  \hline
\end{tabular}
\end{table}

\begin{table}[H]
\centering
\caption{Variación de recaudación tributaria del aumento del PIT (porcentaje del PIB nominal inicial)}
\small
\begin{tabular}{|l|r|r|r|}
  \hline
   & Base (\% del PIB) & Simulado (\% del PIB) & Cambio (p.p. del PIB) \\
  \hline
  Recaudación total & 18.1916 & 18.2442 & 0.0527 \\
  \hline
  IDPC (\textit{FIRMTAX}) & 4.7622 & 4.7277 & -0.0346 \\
  \hline
  PIT & 1.8103 & 1.9798 & 0.1696 \\
  \hline
  IVA & 8.6254 & 8.5633 & -0.0621 \\
  \hline
  Impuestos a la producción & 1.4745 & 1.4640 & -0.0105 \\
  \hline
  Impuestos específicos & 0.9727 & 0.9671 & -0.0056 \\
  \hline
  Aranceles & 0.5465 & 0.5424 & -0.0041 \\
  \hline
\end{tabular}
\end{table}

\begin{table}[H]
\centering
\caption{Descomposición del impacto fiscal del aumento del PIT (puntos del PIB nominal inicial)}
\label{tab:pit-descomposicion}
\small
\resizebox{\textwidth}{!}{%
\begin{tabular}{|l|r|r|r|r|r|}
  \hline
  Impuesto & Cambio observado & Efecto tasa & Efecto base & Tasa más base & Efecto precio \\
  \hline
  IVA & -0.0621 & 0.0000 & -0.0448 & -0.0448 & -0.0173 \\
  \hline
  IDPC (\textit{FIRMTAX}) & -0.0346 & 0.0000 & -0.0288 & -0.0288 & -0.0058 \\
  \hline
  PIT & 0.1696 & 0.1805 & -0.0071 & 0.1734 & -0.0039 \\
  \hline
  Impuestos a la producción & -0.0105 & 0.0000 & -0.0074 & -0.0074 & -0.0032 \\
  \hline
  Impuestos específicos & -0.0056 & 0.0000 & -0.0056 & -0.0056 & 0.0000 \\
  \hline
  Aranceles & -0.0041 & 0.0000 & -0.0031 & -0.0031 & -0.0009 \\
  \hline
  Total & 0.0527 & 0.1805 & -0.0968 & 0.0837 & -0.0310 \\
  \hline
\end{tabular}%
}
\end{table}

El aumento de 10\% en la tasa del impuesto a la renta personal reduce el PIB real en 0.4644\% y el PIB nominal en 0.6742\%. El valor agregado real a costo de factores cae 0.4559\%, por lo que ambas medidas de actividad difieren en menos de 0.01 puntos porcentuales. La contracción del producto es menor que la provocada por el aumento de un punto porcentual del IVA, aunque la caída del consumo es considerable. Este resultado es consistente con un choque que afecta directamente el ingreso disponible de los hogares y, a través de ese canal, reduce tanto el consumo como el ahorro privado.

Por componentes, el consumo real cae 0.6335\%, la inversión real disminuye 0.2541\% y el gasto real del gobierno permanece constante. En el frente externo, las exportaciones reales caen 0.6146\% y las importaciones reales, 0.5683\%, por lo que el saldo comercial disminuye de 2464.05 a 2405.77. En el mercado laboral, la tasa de desempleo aumenta 0.1577 puntos porcentuales para el trabajo de baja calificación, 0.1492 para el de calificación media y 0.1479 para el de alta calificación, con un coeficiente de Okun de 0.326 respecto del valor agregado real. La utilización del capital cae de 82.00\% a 81.50\% de la capacidad instalada ---el capital utilizado se contrae 0.61\%--- y el retorno del capital disminuye 0.12\%.

La variación nominal observada de la recaudación total es de 0.0527 puntos del PIB inicial. Al excluir la revaloración nominal de las bases, el efecto conjunto de tasa y base asciende a 0.0837 puntos del PIB. Este resultado se compone de un efecto tasa mecánico de 0.1805 puntos, asociado íntegramente al PIT, y de un efecto base de -0.0968 puntos, que absorbe el 54\% del rendimiento mecánico. La mayor parte de la erosión proviene de otras bases: el IVA, el IDPC, los impuestos a la producción, los impuestos específicos y los aranceles pierden 0.0448, 0.0288, 0.0074, 0.0056 y 0.0031 puntos, respectivamente, mientras que la base del propio PIT cae 0.0071 puntos. La pérdida de 0.0056 puntos en los impuestos específicos es enteramente de cantidad: refleja el menor consumo físico de combustibles, alcoholes y tabaco, y no admite compensación por precios.

El efecto precio agregado es de -0.0310 puntos del PIB, por lo que reduce la variación observada desde el efecto conjunto de tasa y base de 0.0837 puntos hasta 0.0527 puntos. Sus principales componentes corresponden al IVA y al IDPC, con -0.0173 y -0.0058 puntos, respectivamente.

\textbf{Conclusión del aumento del PIT:} el aumento de 10\% en la tasa genera un efecto conjunto de tasa y base equivalente a 0.0837 puntos del PIB inicial. El efecto mecánico positivo del PIT es atenuado en cerca de la mitad, principalmente por la contracción de las bases del IVA y del impuesto corporativo. Este aumento de la recaudación viene acompañado de una caída del ingreso disponible, el consumo, la inversión, el producto, el empleo y la utilización del capital.

\subsection{Disminución del impuesto a las firmas}

Esta simulación evalúa una disminución de un punto porcentual en la tasa del impuesto de primera categoría (IDPC), equivalente a pasar de 27\% a 26\%\footnote{La carga observada de \textit{FIRMTAX} no coincide con la tasa nominal de referencia de 27\%, porque el agregado incluye tributos empresariales distintos del IDPC. La disminución de un punto porcentual se implementa mediante el factor multiplicativo \(\left(1 - \chi\,\frac{0.01}{0.27}\right)\), con \(\chi=0.916\), participación del IDPC devengado según estadísticas del SII en el impuesto directo pagado por empresas según el Banco Central, para 2023.}.
En el modelo, el agregado \textit{FIRMTAX} incorpora la tributación observada de las firmas en el escenario base; por ello, el choque se interpreta como una reducción de la carga corporativa que afecta directamente la recaudación sobre el capital empresarial.

La Figura~\ref{fig:fit-transmission} sintetiza el mecanismo de transmisión de esta medida. Bajo el cierre con tipo de cambio flexible y ahorro externo fijo, la menor tributación eleva el ingreso neto de las firmas y activa dos canales complementarios: una parte se transmite a los hogares mediante mayores dividendos y consumo, mientras que otra incrementa el ahorro corporativo y la inversión. Como resultado, aumentan la producción, la absorción interna y la utilización de los factores, pero también se intensifica la presión sobre la restricción externa.

\begin{figure}[H]
\centering
\includegraphics[width=\textwidth]{diagrams/diagram_01_disminucion_fit.png}
\caption{Diagrama conceptual de transmisión para la disminución del IDPC.}
\label{fig:fit-transmission}
\end{figure}

\begin{table}[H]
\centering
\caption{Resultados agregados de la disminución del IDPC}
\begin{tabular}{|l|r|r|r|}
  \hline
   & Escenario inicial & Escenario simulado & Cambio \\
  \hline
  PIB nominal & 282097.80 & 282808.52 & 0.2519\% \\
  \hline
  PIB real (gasto) & 282097.80 & 282586.79 & 0.1733\% \\
  \hline
  Valor agregado real (costo de factores) & 249320.66 & 249743.51 & 0.1696\% \\
  \hline
\end{tabular}
\end{table}

\begin{table}[H]
\centering
\caption{Componentes de demanda de la disminución del IDPC}
\begin{tabular}{|l|r|r|}
  \hline
   & Cambio real & Cambio nominal \\
  \hline
  Consumo & 0.1369\% & 0.2129\% \\
  \hline
  Inversión & 0.4116\% & 0.4858\% \\
  \hline
  Gasto de gobierno & 0.0000\% & -- \\
  \hline
  Exportaciones & 0.2034\% & 0.2635\% \\
  \hline
  Importaciones & 0.2425\% & 0.3027\% \\
  \hline
\end{tabular}
\end{table}

\begin{table}[H]
\centering
\caption{Comercio exterior de la disminución del IDPC}
\begin{tabular}{|l|r|r|r|}
  \hline
   & Escenario inicial & Escenario simulado & Cambio \\
  \hline
  Exportaciones & 86815.42 & 87044.20 & 228.78 \\
  \hline
  Importaciones & 84351.36 & 84606.66 & 255.30 \\
  \hline
  Saldo comercial & 2464.05 & 2437.54 & -26.51 \\
  \hline
\end{tabular}
\end{table}

\begin{table}[H]
\centering
\caption{Mercado laboral de la disminución del IDPC}
\begin{tabular}{|l|r|r|r|}
  \hline
   & Escenario inicial & Escenario simulado & Cambio \\
  \hline
  Desempleo, baja calificación & 7.60\% & 7.53\% & -0.0696 p.p. \\
  \hline
  Desempleo, calificación media & 9.90\% & 9.84\% & -0.0587 p.p. \\
  \hline
  Desempleo, alta calificación & 7.80\% & 7.75\% & -0.0533 p.p. \\
  \hline
\end{tabular}
\end{table}

\begin{table}[H]
\centering
\caption{Utilización de factores y respuesta agregada del empleo de la disminución del IDPC}
\small
\begin{tabular}{|l|r|}
  \hline
  Indicador & Valor \\
  \hline
  Utilización del capital, escenario simulado (\% de capacidad instalada; base 82.00\%) & 82.1846\% \\
  \hline
  Capital utilizado ($KD$) & 0.2252\% \\
  \hline
  Retorno del capital ($p^{k}$) & 0.0450\% \\
  \hline
  Desempleo agregado (base 8.3359\%) & 8.2804\% \\
  \hline
  Elasticidad empleo--producto $\eta_{L}$ (respecto del PIB real) & 0.6021 \\
  \hline
  Elasticidad fuerza de trabajo--producto $\eta_{LF}$ (respecto del PIB real) & 0.2529 \\
  \hline
  Coeficiente de Okun respecto del PIB real & 0.3200 \\
  \hline
  Coeficiente de Okun respecto del valor agregado real & 0.3270 \\
  \hline
\end{tabular}
\end{table}

\begin{table}[H]
\centering
\caption{Variación de recaudación tributaria de la disminución del IDPC (porcentaje del PIB nominal inicial)}
\small
\begin{tabular}{|l|r|r|r|}
  \hline
   & Base (\% del PIB) & Simulado (\% del PIB) & Cambio (p.p. del PIB) \\
  \hline
  Recaudación total & 18.1916 & 18.0776 & -0.1140 \\
  \hline
  IDPC (\textit{FIRMTAX}) & 4.7622 & 4.6131 & -0.1491 \\
  \hline
  PIT & 1.8103 & 1.8141 & 0.0039 \\
  \hline
  IVA & 8.6254 & 8.6495 & 0.0242 \\
  \hline
  Impuestos a la producción & 1.4745 & 1.4778 & 0.0033 \\
  \hline
  Impuestos específicos & 0.9727 & 0.9748 & 0.0021 \\
  \hline
  Aranceles & 0.5465 & 0.5482 & 0.0017 \\
  \hline
\end{tabular}
\end{table}

\begin{table}[H]
\centering
\caption{Descomposición del impacto fiscal de la disminución del IDPC (puntos del PIB nominal inicial)}
\label{tab:idpc-descomposicion}
\small
\resizebox{\textwidth}{!}{%
\begin{tabular}{|l|r|r|r|r|r|}
  \hline
  Impuesto & Cambio observado & Efecto tasa & Efecto base & Tasa más base & Efecto precio \\
  \hline
  IVA & 0.0242 & 0.0000 & 0.0177 & 0.0177 & 0.0065 \\
  \hline
  IDPC (\textit{FIRMTAX}) & -0.1491 & -0.1618 & 0.0105 & -0.1512 & 0.0021 \\
  \hline
  PIT & 0.0039 & 0.0000 & 0.0025 & 0.0025 & 0.0014 \\
  \hline
  Impuestos a la producción & 0.0033 & 0.0000 & 0.0021 & 0.0021 & 0.0012 \\
  \hline
  Impuestos específicos & 0.0021 & 0.0000 & 0.0021 & 0.0021 & 0.0000 \\
  \hline
  Aranceles & 0.0017 & 0.0000 & 0.0014 & 0.0014 & 0.0003 \\
  \hline
  Total & -0.1140 & -0.1618 & 0.0363 & -0.1255 & 0.0115 \\
  \hline
\end{tabular}%
}
\end{table}

En términos macroeconómicos, la rebaja del impuesto a las firmas produce una expansión moderada de la actividad: el PIB real aumenta 0.1733\%, el valor agregado real 0.1696\% y el PIB nominal 0.2519\%. Por el lado de la demanda, el consumo real aumenta 0.1369\% y la inversión real, 0.4116\%. Estos resultados son coherentes con los dos canales destacados en la figura: mayores dividendos para los hogares, que elevan el consumo, y un mayor ahorro de las firmas, que impulsa la inversión mediante el bloque ahorro--inversión.

En el frente externo, las exportaciones y las importaciones aumentan 0.2034\% y 0.2425\%, respectivamente, pero el saldo comercial disminuye de 2464.05 a 2437.54. Esta dinámica refleja que la expansión de la absorción interna también incrementa la presión sobre la restricción externa. En el mercado laboral, la tasa de desempleo cae 0.0696 puntos porcentuales para el trabajo de baja calificación, 0.0587 para el de calificación media y 0.0533 para el de alta calificación, con un coeficiente de Okun de 0.327 respecto del valor agregado real. La utilización del capital sube de 82.00\% a 82.18\% de la capacidad instalada, con un aumento de 0.23\% del capital utilizado, y su retorno aumenta 0.05\%.

La variación nominal observada de la recaudación total corresponde a una disminución de 0.1140 puntos del PIB inicial. Al excluir la revaloración nominal de las bases, el efecto conjunto de tasa y base es una caída de 0.1255 puntos del PIB. Este resultado combina un efecto tasa mecánico de -0.1618 puntos, asociado íntegramente al IDPC, y un efecto base positivo de 0.0363 puntos, que compensa cerca de 22\% de la pérdida mecánica. En el propio IDPC, el efecto base es de 0.0105 puntos, reflejo de la mayor utilización del capital corporativo. La expansión económica también eleva las bases del IVA, el PIT, los impuestos a la producción, los impuestos específicos y los aranceles, cuyos aportes ascienden a 0.0177, 0.0025, 0.0021, 0.0021 y 0.0014 puntos del PIB, respectivamente.

El efecto precio agregado es positivo y alcanza 0.0115 puntos del PIB, por lo que reduce la pérdida desde 0.1255 hasta 0.1140 puntos. Su principal componente corresponde a la revaloración de la base del IVA, con 0.0065 puntos del PIB. En el IDPC, el efecto precio de 0.0021 puntos equivale a una quinta parte del efecto base, en línea con $\psi=0.2$. Los impuestos específicos aportan 0.0021 puntos, en su totalidad por cantidad y con efecto precio nulo, como corresponde a un gravamen unitario.

\textbf{Conclusión de la disminución del impuesto a las firmas:} la reducción de un punto porcentual del IDPC genera una pérdida fiscal de tasa más base equivalente a 0.1255 puntos del PIB inicial. La expansión de las bases tributarias ---incluida la del propio IDPC, por la mayor utilización del capital--- compensa parcialmente el efecto mecánico de la rebaja, mientras que el efecto precio positivo reduce la pérdida nominal observada a 0.1140 puntos del PIB. La medida eleva moderadamente el producto, el consumo y, especialmente, la inversión, y mejora el mercado laboral, pero deteriora el saldo comercial. La principal disyuntiva se plantea, por tanto, entre una mayor actividad de corto plazo y menores ingresos fiscales.

\subsection{Aumento de los impuestos específicos sobre productos}
\label{sec:sim-tesp}

Esta simulación evalúa un aumento de 10\% en el impuesto específico por unidad física $\tau_{e}^{p}$ de todos los productos gravados, que en la SAM corresponden esencialmente a combustibles, alcoholes, tabaco y fuentes móviles. El ejercicio solo es posible bajo la especificación vigente: mientras el gravamen se representaba como un monto exógeno de la SAM, su recaudación era constante ante cualquier perturbación y el choque carecía de contenido.

El mecanismo de transmisión difiere del resto de las medidas evaluadas. La cuña actúa sobre el precio de comprador de un conjunto acotado de bienes, por lo que la pérdida de poder adquisitivo se concentra y, al mismo tiempo, se activa una sustitución entre bienes en la función de utilidad CES del hogar. Como el impuesto no grava el uso intermedio, no encarece directamente los costos de producción; el efecto sobre la actividad opera por la vía de la demanda final y del ingreso real.

La Figura~\ref{fig:tesp-transmission} resume este mecanismo. El aumento de $\tau_{e}^{p}$ eleva el precio de comprador del bien gravado y, con ello, la recaudación específica de manera directa. La menor demanda final del bien gravado contrae su producción, el empleo del sector y las bases del IVA, el PIT, los impuestos a la producción y los aranceles, mientras que la base imponible del propio IVA aumenta al incorporar el mayor monto específico. La contracción de la absorción interna en el sector gravado reduce la presión externa y desencadena el ajuste cambiario habitual, con efectos de segundo orden sobre exportaciones e importaciones.

\begin{figure}[H]
\centering
\includegraphics[width=\textwidth]{diagrams/diagram_04_aumento_impuesto_especifico.png}
\caption{Diagrama conceptual de transmisión para el aumento de los impuestos específicos.}
\label{fig:tesp-transmission}
\end{figure}

\begin{table}[H]
\centering
\caption{Resultados agregados del aumento de los impuestos específicos}
\begin{tabular}{|l|r|r|r|}
  \hline
   & Escenario inicial & Escenario simulado & Cambio \\
  \hline
  PIB nominal & 282097.80 & 281512.75 & -0.2074\% \\
  \hline
  PIB real (gasto) & 282097.80 & 281452.10 & -0.2289\% \\
  \hline
  Valor agregado real (costo de factores) & 249320.66 & 248780.63 & -0.2166\% \\
  \hline
\end{tabular}
\end{table}

\begin{table}[H]
\centering
\caption{Componentes de demanda del aumento de los impuestos específicos}
\begin{tabular}{|l|r|r|}
  \hline
   & Cambio real & Cambio nominal \\
  \hline
  Consumo & -0.3037\% & -0.2517\% \\
  \hline
  Inversión & -0.1615\% & -0.1362\% \\
  \hline
  Gasto de gobierno & 0.0000\% & -- \\
  \hline
  Exportaciones & -0.4262\% & -0.5006\% \\
  \hline
  Importaciones & -0.4193\% & -0.4936\% \\
  \hline
\end{tabular}
\end{table}

\begin{table}[H]
\centering
\caption{Comercio exterior del aumento de los impuestos específicos}
\begin{tabular}{|l|r|r|r|}
  \hline
   & Escenario inicial & Escenario simulado & Cambio \\
  \hline
  Exportaciones & 86815.42 & 86380.83 & -434.59 \\
  \hline
  Importaciones & 84351.36 & 83935.00 & -416.36 \\
  \hline
  Saldo comercial & 2464.05 & 2445.82 & -18.23 \\
  \hline
\end{tabular}
\end{table}

\begin{table}[H]
\centering
\caption{Mercado laboral del aumento de los impuestos específicos}
\begin{tabular}{|l|r|r|r|}
  \hline
   & Escenario inicial & Escenario simulado & Cambio \\
  \hline
  Desempleo, baja calificación & 7.60\% & 7.67\% & 0.0747 p.p. \\
  \hline
  Desempleo, calificación media & 9.90\% & 9.96\% & 0.0630 p.p. \\
  \hline
  Desempleo, alta calificación & 7.80\% & 7.86\% & 0.0576 p.p. \\
  \hline
\end{tabular}
\end{table}

\begin{table}[H]
\centering
\caption{Utilización de factores y respuesta agregada del empleo del aumento de los impuestos específicos}
\small
\begin{tabular}{|l|r|}
  \hline
  Indicador & Valor \\
  \hline
  Utilización del capital, escenario simulado (\% de capacidad instalada; base 82.00\%) & 81.7499\% \\
  \hline
  Capital utilizado ($KD$) & -0.3050\% \\
  \hline
  Retorno del capital ($p^{k}$) & -0.0611\% \\
  \hline
  Desempleo agregado (base 8.3359\%) & 8.3957\% \\
  \hline
  Elasticidad empleo--producto $\eta_{L}$ (respecto del PIB real) & 0.4918 \\
  \hline
  Elasticidad fuerza de trabajo--producto $\eta_{LF}$ (respecto del PIB real) & 0.2067 \\
  \hline
  Coeficiente de Okun respecto del PIB real & 0.2614 \\
  \hline
  Coeficiente de Okun respecto del valor agregado real & 0.2763 \\
  \hline
\end{tabular}
\end{table}

\begin{table}[H]
\centering
\caption{Variación de recaudación tributaria del aumento de los impuestos específicos (porcentaje del PIB nominal inicial)}
\small
\begin{tabular}{|l|r|r|r|}
  \hline
   & Base (\% del PIB) & Simulado (\% del PIB) & Cambio (p.p. del PIB) \\
  \hline
  Recaudación total & 18.1916 & 18.2357 & 0.0441 \\
  \hline
  IDPC (\textit{FIRMTAX}) & 4.7622 & 4.7448 & -0.0174 \\
  \hline
  PIT & 1.8103 & 1.8057 & -0.0046 \\
  \hline
  IVA & 8.6254 & 8.6061 & -0.0193 \\
  \hline
  Impuestos a la producción & 1.4745 & 1.4707 & -0.0039 \\
  \hline
  Impuestos específicos & 0.9727 & 1.0648 & 0.0921 \\
  \hline
  Aranceles & 0.5465 & 0.5436 & -0.0029 \\
  \hline
\end{tabular}
\end{table}

\begin{table}[H]
\centering
\caption{Descomposición del impacto fiscal del aumento de los impuestos específicos (puntos del PIB nominal inicial)}
\label{tab:tesp-descomposicion}
\small
\resizebox{\textwidth}{!}{%
\begin{tabular}{|l|r|r|r|r|r|}
  \hline
  Impuesto & Cambio observado & Efecto tasa & Efecto base & Tasa más base & Efecto precio \\
  \hline
  IVA & -0.0193 & 0.0000 & -0.0277 & -0.0277 & 0.0084 \\
  \hline
  IDPC (\textit{FIRMTAX}) & -0.0174 & 0.0000 & -0.0145 & -0.0145 & -0.0029 \\
  \hline
  PIT & -0.0046 & 0.0000 & -0.0030 & -0.0030 & -0.0015 \\
  \hline
  Impuestos a la producción & -0.0039 & 0.0000 & -0.0025 & -0.0025 & -0.0014 \\
  \hline
  Impuestos específicos & 0.0921 & 0.0970 & -0.0049 & 0.0921 & 0.0000 \\
  \hline
  Aranceles & -0.0029 & 0.0000 & -0.0025 & -0.0025 & -0.0004 \\
  \hline
  Total & 0.0441 & 0.0970 & -0.0552 & 0.0419 & 0.0022 \\
  \hline
\end{tabular}%
}
\end{table}

El alza de 10\% en los gravamenes específicos reduce el valor agregado real en 0.2166\% y el PIB real por el gasto en 0.2289\%. El PIB nominal, en cambio, cae solo 0.2074\%: la diferencia refleja que el PIB a precios de mercado incorpora los impuestos sobre productos, de modo que la mayor cuña sostiene el valor nominal mientras los volúmenes se contraen. El mismo patrón aparece en el consumo, cuya caída real de 0.3037\% supera su caída nominal de 0.2517\%. Esta separación entre magnitudes reales y nominales solo es observable porque los impuestos indirectos operan ahora como cuñas de precio.

En los mercados de factores, la tasa de desempleo aumenta entre 0.0576 y 0.0747 puntos porcentuales según el tipo de trabajo, y la utilización del capital cae de 82.00\% a 81.75\% de la capacidad instalada. El coeficiente de Okun respecto del valor agregado real es 0.276, por debajo de los valores de 0.326 y 0.327 obtenidos en los choques de PIT e IDPC y del extremo inferior de la banda empírica de referencia. La explicación es de composición: la medida concentra su impacto en un subconjunto de productos cuyos sectores proveedores tienen una intensidad laboral menor que el promedio de la economía, de modo que una misma contracción del valor agregado se traduce en menos empleo perdido. El resultado ilustra una limitación de la aproximación agregada usada en la calibración de $\phi$, $\psi$ y $\varepsilon$, que no distingue el origen sectorial del choque.

Desde el punto de vista recaudatorio, el rendimiento mecánico asciende a 0.0970 puntos del PIB y la erosión de la propia base del impuesto es de apenas 0.0049 puntos, esto es, 5\% del efecto tasa. Se trata de la relación tasa--base más favorable de las cuatro simulaciones y es coherente con la baja elasticidad precio que la especificación CES atribuye a estos bienes en el nivel de agregación del modelo. La atenuación relevante proviene de las demás bases: el IVA pierde 0.0277 puntos, el IDPC 0.0145, el PIT 0.0030, los impuestos a la producción 0.0025 y los aranceles 0.0025, lo que suma 0.0503 puntos. El efecto conjunto de tasa y base sobre la recaudación total es, por tanto, de 0.0419 puntos del PIB, menos de la mitad del rendimiento mecánico.

El efecto precio agregado es de 0.0022 puntos del PIB y resulta de la compensación entre una revaloración positiva de la base del IVA ---que se calcula sobre el precio con impuesto específico incluido--- y efectos precio negativos en los demás instrumentos. En la fila de impuestos específicos, el efecto precio es exactamente nulo: la descomposición bilineal de la Sección~\ref{sec:metodologia-descomp} se cumple sin residuo y confirma numéricamente que el gravamen unitario no se revaloriza con los precios.

\textbf{Conclusión del aumento de los impuestos específicos:} un alza de 10\% en los montos por unidad física genera un efecto conjunto de tasa y base de 0.0419 puntos del PIB inicial, con un costo en actividad de 0.2166\% del valor agregado real. La relación entre costo en actividad y recaudación de tasa más base es de 5.2 puntos de valor agregado por punto del PIB recaudado, similar a la del aumento del PIT (5.4) pero concentrada en un conjunto acotado de bienes, lo que tiene implicancias distributivas que este modelo de hogar representativo no permite evaluar y que requieren la extensión por quintiles.

\section{Conclusión general}

El documento desarrolla un modelo de equilibrio general computable estático y determinístico para una economía de doce sectores y trece bienes, calibrado con una matriz de contabilidad social para Chile. El modelo integra los bloques de hogares, producción, firmas, gobierno, sector externo, recaudación tributaria y ahorro--inversión. También articula la oferta doméstica, las exportaciones y las importaciones mediante los mecanismos CET y Armington. En los mercados de factores, la utilización endógena del capital y la participación laboral endógena permiten que el ajuste de las cantidades se distribuya entre capital y trabajo, y que la respuesta del desempleo al producto se mantenga dentro del rango observado para Chile.

Los ejercicios muestran que las cuatro reformas tributarias producen efectos agregados diferenciados. El aumento de 10\% del impuesto a la renta de personas reduce el valor agregado real en 0.4559\%, el consumo real en 0.6335\% y la inversión real en 0.2541\%, e incrementa la tasa de desempleo entre 0.1479 y 0.1577 puntos porcentuales. El alza de 10\% de los impuestos específicos reduce el valor agregado real en 0.2166\%, el consumo real en 0.3037\% y la inversión real en 0.1615\%, y eleva el desempleo entre 0.0576 y 0.0747 puntos. La disminución de un punto porcentual del IDPC eleva el PIB real en 0.1733\% y la inversión real en 0.4116\%, y reduce la tasa de desempleo entre 0.0533 y 0.0696 puntos porcentuales, aunque deteriora el saldo comercial. El incremento de un punto porcentual en la tasa del IVA reduce el valor agregado real en 0.9105\%, el consumo real en 1.2603\% y la inversión real en 0.6549\%, y aumenta la tasa de desempleo entre 0.2879 y 0.3313 puntos porcentuales.

Desde el punto de vista recaudatorio, al considerar el efecto conjunto de tasa y base, el incremento del IVA eleva la recaudación total en 0.2371 puntos del PIB inicial, el aumento del PIT en 0.0837 puntos y el alza de los impuestos específicos en 0.0419 puntos, mientras que la disminución del IDPC la reduce en 0.1255 puntos. Las variaciones nominales observadas son de 0.1755, 0.0527, 0.0441 y -0.1140 puntos, respectivamente, debido a efectos precio de -0.0616, -0.0310, 0.0022 y 0.0115 puntos del PIB. En los aumentos tributarios, la contracción de las bases absorbe entre la mitad y algo más de la mitad del rendimiento mecánico; en la rebaja del IDPC, la expansión de las bases compensa cerca de 22\% de la pérdida. Una parte relevante de estos efectos base proviene del impuesto corporativo, cuya base responde al grado de utilización del capital.

Ordenados por costo en actividad por punto del PIB recaudado, el IVA resulta el instrumento menos oneroso (3.8 puntos de valor agregado real), seguido de los impuestos específicos (5.2) y del impuesto a la renta de personas (5.4). Esta jerarquía depende de manera decisiva de que el IVA otorgue crédito pleno sobre el uso intermedio: bajo la especificación anterior, que gravaba implícitamente los insumos, el mismo choque aparecía como el más costoso de los tres. La comparación entre instrumentos revela además un patrón común: la erosión de las bases distintas de la gravada explica una parte sustantiva de la atenuación del rendimiento mecánico. En el alza de los impuestos específicos, la propia base cede solo 5\% del efecto tasa, pero la contracción de las restantes bases resta 0.0503 puntos del PIB, más de la mitad del rendimiento. Este resultado refuerza el argumento a favor de la evaluación en equilibrio general: una estimación de equilibrio parcial del mismo choque, centrada en la elasticidad de la demanda del bien gravado, habría proyectado un rendimiento cercano al mecánico.

Una ventaja de este tipo de modelos es que permiten evaluar simultáneamente los efectos de una medida sobre variables reales, externas, laborales y fiscales, y no solo sobre la dimensión que motiva originalmente la reforma. A diferencia de los análisis parciales, esta perspectiva integrada hace visibles las disyuntivas entre actividad, empleo, recaudación y sector externo, y ofrece una base más completa para la evaluación de la política tributaria.

Los resultados deben interpretarse considerando algunas limitaciones. Las dos primeras observaciones de versiones anteriores ---el registro del IVA como volumen en el PIB real por el gasto y la ausencia de distinción entre precio de comprador y precio de vendedor--- quedan resueltas por la especificación descrita en la Sección~\ref{sec:productos}, que separa $pq^{p}$ de $p^{p}$ y deja el vaciado de mercado en unidades físicas.

Permanecen, sin embargo, otras limitaciones. Primero, los parámetros de los mercados de factores se calibran para reproducir un coeficiente de Okun obtenido de una única fuente, y la elasticidad de la oferta de servicios de capital carece de un ancla empírica directa para Chile. Segundo, la calibración supone que el consumo intermedio no paga impuestos específicos, supuesto que subestima la carga sobre los sectores intensivos en combustibles ---transporte y minería, en particular--- y que la evidencia sobre recuperación parcial del impuesto al diésel permitiría refinar. Tercero, el residuo $t_{iva,int}^{p}$ representa el IVA no recuperable de manera agregada y transitoria, en lugar de derivarlo de la estructura de exenciones y de la información del formulario 29. Cuarto, el traspaso del impuesto a precios es unitario por construcción, supuesto razonable bajo competencia perfecta pero que la literatura sobre incidencia de impuestos selectivos matiza en presencia de poder de mercado \cite{delipalla_keen_1992}. Quinto, el nivel base de utilización del capital renormaliza los resultados sin imponer un techo de capacidad; incorporar la restricción $u_{k}\leq 1$ es la extensión natural de este margen.

Como próximas etapas, se propone estimar directamente las elasticidades del empleo y de la fuerza de trabajo respecto del producto con datos de la Encuesta Nacional de Empleo, e implementar un análisis de sensibilidad que incluya los parámetros $\phi$, $\psi$ y $\varepsilon$. Entre las extensiones posibles se encuentran una representación más detallada de la estructura tributaria, una mayor desagregación sectorial y de productos, la incorporación de hogares heterogéneos, el análisis distributivo y una profundización del mercado laboral mediante rigideces adicionales o mecanismos de informalidad.


\section{Disponibilidad de código}

El repositorio del proyecto se encuentra disponible en \url{https://github.com/fcohenriquez/eq_general}.

\renewcommand{\refname}{Referencias}
\nocite{*}
\bibliographystyle{plain}
\bibliography{refs}

\end{document}